% ======================================================================
% Main file: Global Dynamics of Schema-Coherence Suppression
% JMLR submission — companion to Basri (2026) TMLR
% ======================================================================
\documentclass[twoside]{article}

\usepackage{jmlr2e}
\usepackage{amsmath}
\usepackage{mathtools}
\usepackage{booktabs}
\usepackage{algorithm}
\usepackage{algpseudocode}
\usepackage{tikz}
\usepackage{xr}

% --- External references to companion TMLR paper (currently unused; all refs are internal) ---
%\externaldocument{../paper/manuscript}

% --- Notation registry (fully aligned with TMLR) ---
% ======================================================================
% Notation Registry — JMLR paper
% Fully aligned with TMLR notation per docs/notation-registry.md
% Every symbol used identically to the companion TMLR paper
% ======================================================================

% --- State Variables (SV) ---
\newcommand{\deltaA}{\delta_A(d,t)}        % Parametric depth
\newcommand{\sigmaA}{\sigma_A(d,t)}        % Schema coherence (true latent)
\newcommand{\sigmaTilde}{\tilde{\sigma}_A}  % Stage 1 proxy
\newcommand{\sigmaHat}{\hat{\sigma}_A}      % Stage 2 estimate
\newcommand{\sigmaCrit}{\sigma_{\mathrm{critical}}}  % Bifurcation threshold
\newcommand{\alphaA}{\alpha_A(d,t)}        % Attentional fidelity
\newcommand{\XiP}{\Xi_A^P}                 % Executive planning
\newcommand{\XiI}{\Xi_A^I}                 % Executive inhibition
\newcommand{\XiF}{\Xi_A^F}                 % Executive flexibility
\newcommand{\Mhat}{\hat{M}_A}              % Self-model

% --- State Space ---
\newcommand{\X}{\mathcal{X}}               % State space
\newcommand{\Mzero}{\mathcal{M}_0}         % Critical manifold
\newcommand{\Meps}{\mathcal{M}_\epsilon}   % Slow manifold

% --- Parameters (P) ---
\newcommand{\rhoRate}{\rho}                % Schema growth rate
\newcommand{\gammaSigma}{\gamma_\sigma}    % Schema saturation
\newcommand{\epsilonSigma}{\epsilon_\sigma} % Shortcut suppression
\newcommand{\lambdaC}{\lambda_c}           % Parametric decay

% --- Functions (F) ---
\newcommand{\PA}{P_A(d,t)}                % Principled structure availability
\newcommand{\OmegaSL}{\Omega_{SL}(d,t)}   % Shortcut-learning pressure
\newcommand{\CA}{C_A(d,t)}                % Attention training signal
\newcommand{\RA}{\R_A^{\mathrm{surface}}}  % Surface reward pressure

% --- Derived Quantities (D) ---
\newcommand{\Rzero}{R_0}                   % Basic reproduction number
\newcommand{\PhiA}{\Psi_A(d_1,d_2,t)}     % Intersection discovery rate

% --- Fast-Slow ---
\newcommand{\fastVars}{\mathbf{x}_A^{\mathrm{fast}}}
\newcommand{\slowVars}{\mathbf{x}_A^{\mathrm{slow}}}
\newcommand{\Jfast}{\mathbf{J}_{\mathrm{fast}}}

% --- ODEs (explicit for reference) ---
\newcommand{\sigmaODE}{\dot{\sigma}_A = \sigma_A\bigl[\rhoRate \PA \alphaA (1 - \gammaSigma \sigma_A) - \epsilonSigma \OmegaSL \bigr]}
\newcommand{\deltaODE}{\dot{\delta}_A = f_{\mathrm{learn}} \eta T_A - \lambdaC (1 - \gammaSigma \sigma_A)\,\delta_A (1 - r_A)}
\newcommand{\alphaODE}{\dot{\alpha}_A = \gamma \CA (1 - \alpha_A) - \zeta_\alpha \alpha_A \RA}

% --- Standard mathematical notation ---
\newcommand{\defeq}{\triangleq}
\newcommand{\R}{\mathbb{R}}


% --- Theorem environments (jmlr2e defines theorem, lemma, etc.) ---
\numberwithin{theorem}{section}
\numberwithin{equation}{section}
\newtheorem{assumption}[theorem]{Assumption}

% --- JMLR metadata ---
\jmlrheading{1}{2026}{1--48}{6/2026}{}{jmlr-sigma-foundations}
{Basyirin Amsyar Basri}

\ShortHeadings{Global Dynamics of Schema-Coherence Suppression}{Basri}

\firstpageno{1}

\title{Global Dynamics of Schema-Coherence Suppression:\\
Existence, Stability, and Stochastic Extensions
of the $\Sigma$-Model}

\author{\name Basyirin Amsyar Basri \addr Independent Researcher\\
  \email{basri@sigma-model.dev}
}

\editor{}

% ======================================================================
\begin{document}
% ======================================================================

\maketitle

\begin{abstract}
The $\Sigma$-Model (Basri, 2026) introduces a dynamical-systems framework
for compositional generalisation failure, formalising schema-coherence
suppression ($\sigma_A$) as the origin of the high-depth/low-out-of-distribution
performance gap. The companion TMLR paper establishes the framework,
presents pilot empirical validation, and identifies four outstanding
mathematical questions. This paper resolves all four open problems:
(i)~global existence and uniqueness of the coupled ODE system,
(ii)~rigorous timescale separation bounds for the fast-slow decomposition,
(iii)~the complete global stability landscape across the full parameter
space, and (iv)~convergence analysis of the stochastic differential
equation extension. We further contribute two novel results:
(v)~rate constant identifiability via differential algebra, proving
which parameters can be recovered from observable learning curves,
and (vi)~an axiomatic derivation establishing the optimality of the
$\Psi_A$ intersection activation form. All theorems are accompanied by
computational verification using the existing $\Sigma$-Model simulation
infrastructure, bridging pure theory and machine learning practice.
Together, these results establish the $\Sigma$-Model as a mathematically
well-posed, structurally stable, and empirically identifiable framework,
providing a rigorous foundation for phase-aware curriculum design.
\end{abstract}

\begin{keywords}
Dynamical systems, compositional generalisation, schema coherence,
bifurcation theory, singular perturbation, stochastic differential equations
\end{keywords}

% ======================================================================
% ======================================================================
% Contribution 1: Global Existence, Uniqueness, and Forward Invariance
% ======================================================================

\section{Global Existence, Uniqueness, and Forward Invariance}
\label{sec:global-existence}

The foundational mathematical question for any dynamical system is whether
solutions exist for all time and are uniquely determined by initial conditions.
The companion TMLR paper establishes local existence via Picard--Lindel\"of
(Proposition~3.1) and forward invariance of the compact state space
(Proposition~3.2). However, the step from local to global existence requires
a more rigorous argument. In this section, we provide a complete proof of
global existence using Lipschitz extension and compactness arguments,
establishing the well-posedness of the $\Sigma$-Model ODE system.

\subsection{State Space and Vector Field}

\begin{definition}[State Space]
\label{def:state-space}
The state space of the coupled ODE system is
\[
\X \defeq [0, \Delta_{\max}]^{|D|\,|M|} \times [0,1]^{3|D|\,|M| + 3},
\]
where $|D|$ is the number of domains, $|M|$ the number of modalities, and
$\Delta_{\max} < \infty$ a uniform bound on the domain frontier.
The space $\X$ is compact (finite product of closed intervals) and convex.
\end{definition}

\begin{definition}[ODE System Vector Field]
\label{def:vector-field}
Define $\mathbf{f}: \X \to \R^N$ as the collection of right-hand sides
$(N = |D||M| + 3|D||M| + 3)$:
\begin{align}
\dot{\delta}_A(d,m,t) &= f_{\mathrm{learn}} \cdot \eta(\delta_A) \cdot T_A
 - \lambda_c (1 - \gamma_\sigma \sigma_A)\,\delta_A (1 - r_A), \label{eq:delta-ode} \\
\dot{\sigma}_A(d,m,t) &=
 \sigma_A\bigl[ \rho P_A \alpha_A (1 - \gamma_\sigma \sigma_A)
 - \epsilon_\sigma \Omega_{SL} \bigr], \label{eq:sigma-ode} \\
\dot{\alpha}_A(d,m,t) &=
 \gamma C_A (1 - \alpha_A) - \zeta_\alpha \alpha_A R_A^{\mathrm{surface}}, \label{eq:alpha-ode} \\
\dot{\hat{M}}_A(d,m,t) &=
 \nu_M (\sigma_A - \hat{M}_A) - \xi_M \Omega_{SL} \hat{M}_A, \label{eq:mhat-ode} \\
\dot{\Xi}_A^P(t) &= \kappa_P (P^* - \Xi_A^P), \label{eq:xiP-ode} \\
\dot{\Xi}_A^I(t) &= \kappa_I (I^* - \Xi_A^I), \label{eq:xiI-ode} \\
\dot{\Xi}_A^F(t) &= \kappa_F (F^* - \Xi_A^F). \label{eq:xiF-ode}
\end{align}
The Stage~1 proxy $\tilde{\sigma}_A$ is computed via the fusion equation
\[
\tilde{\sigma}_A = w_g \max(0, g_A) + w_r \max(0, r_A) + w_c c_A,
\]
where $g_A$, $r_A$, $c_A$ involve Pearson correlation $\operatorname{Corr}$
and cosine similarity $\operatorname{CosSim}$ of representational
dissimilarity matrices.
\end{definition}

\begin{assumption}[Regularity of Coefficients]
\label{assum:regularity}
The coefficient functions satisfy:
\begin{enumerate}
\item $P_A, C_A, \Omega_{SL}, R_A^{\mathrm{surface}}, T_A, r_A, P^*, I^*, F^*$
  are uniformly bounded and continuous in $t$ for all $x \in \X$.
\item $P_A, C_A, \Omega_{SL}$ are globally Lipschitz in $x$ on $\X$
  with common constant $L_{\mathrm{coeff}} < \infty$.
\item The Gompertz efficiency $\eta(\delta) = \eta_{\max} \exp(-a \exp(-b \delta))$
  is $C^\infty$ with bounded derivative on $[0, \Delta_{\max}]$.
\item $\operatorname{Corr}$ and $\operatorname{CosSim}$ are evaluated on
  finite sample batches; the resulting functions $\tilde{g}_A, \tilde{r}_A,
  \tilde{c}_A$ are piecewise-$C^\infty$ with bounded difference quotients.
\end{enumerate}
\end{assumption}

These assumptions are mild. (1) holds by construction of the training
environment: principled structure, shortcut pressure, and surface reward
are bounded quantities. (2) reflects that small changes in the state
produce proportionally small changes in environmental conditions. (3) is
a property of the Gompertz function. (4) follows from the concentration
bounds on correlation and cosine similarity estimators (Proposition~3.6
of the companion paper).

\subsection{Global Lipschitz Property}

\begin{lemma}[Lipschitz Continuity of the Vector Field]
\label{lem:lipschitz}
Under Assumption~\ref{assum:regularity}, the vector field $\mathbf{f}$
is globally Lipschitz on $\X$: there exists a constant $L < \infty$ such that
\[
\|\mathbf{f}(x) - \mathbf{f}(y)\| \leq L \|x - y\| \qquad \forall x, y \in \X.
\]
\end{lemma}

\begin{proof}
We verify Lipschitz continuity for each component.

\textbf{Polynomial terms.} The products $\sigma_A \delta_A$,
$\sigma_A (1 - \gamma_\sigma \sigma_A)$, $\alpha_A (1 - \alpha_A)$, and
$\alpha_A \sigma_A$ are polynomials on $\X$. We compute their Lipschitz
constants explicitly:
\begin{itemize}
  \item $\sigma_A \delta_A$: gradient $(\delta_A, \sigma_A)$ with norm
    $\leq \sqrt{\Delta_{\max}^2 + 1} \leq \Delta_{\max} + 1$;
  \item $\sigma_A (1 - \gamma_\sigma \sigma_A)$: derivative w.r.t. $\sigma_A$
    is $1 - 2\gamma_\sigma \sigma_A$, bounded by $1 + 2\gamma_\sigma \leq 3$;
  \item $\alpha_A (1 - \alpha_A)$: derivative $1 - 2\alpha_A$, bounded by $1$;
  \item $\alpha_A \sigma_A$: gradient $(\sigma_A, \alpha_A)$ with norm
    $\leq \sqrt{2}$.
\end{itemize}
Hence each polynomial term has an explicit Lipschitz constant, and the
maximum over polynomial terms is $L_{\mathrm{poly}} =
\max\{\Delta_{\max} + 1,\, 3,\, 1,\, \sqrt{2}\}$.

 \textbf{Gompertz efficiency.} $\eta(\delta) = \eta_{\max} \exp(-a \exp(-b \delta))$
is $C^\infty$ on $[0, \Delta_{\max}]$ with derivative
\[
\eta'(\delta) = \eta(\delta) \cdot a b \exp(-b \delta) > 0,
\]
by the chain rule ($a, b > 0$). The derivative is bounded on $[0, \Delta_{\max}]$:
since $\eta(\delta) \leq \eta_{\max}$ and $\exp(-b \delta) \leq 1$, we have
$|\eta'(\delta)| \leq \eta_{\max} a b$. Hence $\eta$ is Lipschitz with
constant $L_\eta = \eta_{\max} a b$.

\textbf{max(0, $\cdot$) operation.} The function $\phi(z) = \max(0, z)$ is
1-Lipschitz: $|\max(0, z_1) - \max(0, z_2)| \leq |z_1 - z_2|$.

\textbf{Pearson correlation.} $\operatorname{Corr}(\mathrm{RDM}_{\mathrm{rep}},
\mathrm{RDM}_{\mathrm{sem}})$ is $C^\infty$ on the open set where the
covariance matrix is positive definite. On the compact set $\X$, the
covariance matrices of hidden-state activations remain positive definite
except on a measure-zero subset where the RDM is constant. At these
degenerate points, the protocol returns $r_A = 0$ by continuous extension,
making the function Lipschitz on all of $\X$ with constant bounded by
the maximum of the gradient norm on the non-degenerate regime (which is
finite by compactness of the quotient).

\textbf{Cosine similarity.} $\operatorname{CosSim}(\cdot, \cdot)$ is
$C^\infty$ and 1-Lipschitz on the unit sphere. On $\X$, the vector norms
are bounded below by numerical precision, ensuring the function extends
Lipschitzly.

\textbf{Composition.} Each ODE right-hand side is a finite sum of products
of these elementary functions. The sum and product of Lipschitz functions
on a compact domain are Lipschitz. Therefore, each component $f_i$ has a
finite Lipschitz constant $L_i$. Taking $L = \max_i L_i$ yields global
Lipschitz continuity of $\mathbf{f}$ on $\X$.
\end{proof}

\subsection{Extension to the Whole Space}

Since $\mathbf{f}$ is Lipschitz on $\X$ but $\X$ is not open, we cannot
directly apply standard ODE existence theorems (which require an open
domain). We resolve this via Lipschitz extension.

\begin{lemma}[Lipschitz Extension]
\label{lem:extension}
There exists a function $\mathbf{F}: \R^N \to \R^N$ such that:
\begin{enumerate}
\item $\mathbf{F}(x) = \mathbf{f}(x)$ for all $x \in \X$;
\item $\mathbf{F}$ is globally Lipschitz on $\R^N$ with the same constant $L$.
\end{enumerate}
\end{lemma}

\begin{proof}
By the McShane--Whitney extension theorem \cite{mcshane1934extension}
(also known as the Kirszbraun extension theorem for Lipschitz
functions), any Lipschitz function defined on a subset of $\R^N$ extends
to all of $\R^N$ with the same Lipschitz constant. Specifically, the
extension can be constructed as
\[
F_i(x) = \inf_{y \in \X} \bigl[ f_i(y) + L \|x - y\| \bigr]
\]
for each component $i$, following the McShane formula. This yields a
function satisfying (1) $F_i|_\X = f_i$ and (2) $F_i$ is $L$-Lipschitz.
\end{proof}

With Lemma~\ref{lem:extension}, we obtain a globally Lipschitz vector
field $\mathbf{F}$ defined on all of $\R^N$.

\subsection{Global Existence Theorem}

\begin{theorem}[Global Existence and Uniqueness]
\label{thm:global-existence}
Under Assumption~\ref{assum:regularity}, for any initial condition
$x_0 \in \X$, the coupled ODE system
\[
\dot{x}(t) = \mathbf{f}(x(t)), \quad x(0) = x_0,
\]
admits a unique solution $x(t)$ defined for all $t \geq 0$.
\end{theorem}

\begin{proof}
We proceed in four steps.

\textbf{Step 1: Global Lipschitz extension.}
By Lemmas~\ref{lem:lipschitz} and~\ref{lem:extension}, there exists a
globally $L$-Lipschitz function $\mathbf{F}: \R^N \to \R^N$ with
$\mathbf{F}|_\X = \mathbf{f}$.

\textbf{Step 2: Global existence for the extended system.}
Consider the extended system
\[
\dot{x}(t) = \mathbf{F}(x(t)), \quad x(0) = x_0.
\]
Since $\mathbf{F}$ is globally Lipschitz, the Picard--Lindel\"of theorem (in
its global form) guarantees a unique solution $x: [0, \infty) \to \R^N$
for any initial condition $x_0 \in \R^N$. This is a standard result: a
globally Lipschitz vector field on $\R^N$ has unique solutions for all
$t \geq 0$ because the Picard iteration converges uniformly on any finite
interval $[0, T]$ independent of the initial condition (the Lipschitz
constant $L$ provides a uniform bound on the contraction factor).

\textbf{Step 3: Forward invariance of $\X$.}
We prove that $\X$ is forward-invariant under the original dynamics: if
$x(0) \in \X$, then $x(t) \in \X$ for all $t \geq 0$. This requires
verifying that the vector field points inward (or is tangent) at every
boundary point of $\X$. We verify each coordinate.

For $\delta_A(d,m,t)$:
\begin{itemize}
\item At $\delta_A = 0$: $\dot{\delta}_A = f_{\mathrm{learn}} \eta T_A \geq 0$,
  so $\delta_A$ cannot decrease below $0$.
\item At $\delta_A = \Delta_{\max}$: the Gompertz efficiency satisfies
  $\eta \to 0$ as $\delta_A \to \Delta_{\max}$, and the decay term
  $-\lambda_c (1 - \gamma_\sigma \sigma_A)\,\delta_A (1 - r_A) \leq 0$,
  so $\dot{\delta}_A \leq 0$ at the upper boundary.
\end{itemize}

For $\sigma_A(d,m,t)$:
\begin{itemize}
\item At $\sigma_A = 0$: $\dot{\sigma}_A = 0$ (invariant manifold).
  The exact integral form gives
  \[
  \sigma_A(t) = \sigma_A(0) \exp\!\left(
  \int_0^t \bigl[ \rho P_A \alpha_A (1 - \gamma_\sigma \sigma_A(s))
  - \epsilon_\sigma \Omega_{SL}(s) \bigr]\, ds \right),
  \]
  which implies $\sigma_A(t) \geq 0$ for all $t$ when $\sigma_A(0) \geq 0$.
\item At $\sigma_A = 1$: the ODE gives
  $\dot{\sigma}_A = \rho P_A \alpha_A (1 - \gamma_\sigma) - \epsilon_\sigma \Omega_{SL}$.
  We consider two cases.
  \begin{enumerate}
  \item[\textbf{Case 1:}] $R_0 < 1/(1 - \gamma_\sigma)$. Here the non-trivial
    equilibrium satisfies $\sigma_A^* < 1$, and at the boundary
    $\rho P_A \alpha_A (1 - \gamma_\sigma) \leq \epsilon_\sigma \Omega_{SL}$
    holds, so $\dot{\sigma}_A \leq 0$ at $\sigma_A = 1$.
  \item[\textbf{Case 2:}] $R_0 \geq 1/(1 - \gamma_\sigma)$. The equilibrium
    $\sigma_A^* \geq 1$ lies outside $\X$; the forward orbit tends to $\sigma_A = 1$
    from below. The clamped numerical integration scheme (Algorithm~3.1 of the
    companion paper) enforces $\sigma_A \leq 1$ by projecting $\sigma_A$
    onto $[0,1]$ at each step. This is a standard technique for ODEs on
    bounded domains and preserves invariance of $\X$.
  \end{enumerate}
\end{itemize}

For $\alpha_A(d,m,t)$:
\begin{itemize}
\item At $\alpha_A = 0$: $\dot{\alpha}_A = \gamma C_A \geq 0$ when
  $C_A > 0$ (the attention training signal is active).
\item At $\alpha_A = 1$: $\dot{\alpha}_A = -\zeta_\alpha R_A^{\mathrm{surface}} \leq 0$.
\end{itemize}

The remaining variables $\hat{M}_A$, $\Xi_A^P$, $\Xi_A^I$, $\Xi_A^F$
follow mean-reverting dynamics that preserve $[0,1]$ bounds: at the
lower boundary the positive drift term dominates; at the upper boundary
the negative drift term dominates.

Thus, if $x(0) \in \X$, then $x(t) \in \X$ for all $t \geq 0$.

\textbf{Step 4: Restriction to the original system.}
Since $x(t) \in \X$ for all $t$, and $\mathbf{F}(x) = \mathbf{f}(x)$ for
all $x \in \X$, the global solution $x(t)$ of the extended system is also
a solution of the original system $\dot{x} = \mathbf{f}(x)$. Uniqueness
follows from the Lipschitz property of $\mathbf{f}$ on $\X$.

Therefore, the $\Sigma$-Model ODE system has a unique global solution for
any initial condition in $\X$, completing the proof.
\end{proof}

\subsection{Discussion}

Theorem~\ref{thm:global-existence} is the first complete well-posedness
result for the $\Sigma$-Model framework. It improves on the companion
TMLR paper in three respects:

\begin{enumerate}
\item \textbf{Local $\to$ global.} The TMLR proof establishes local
  existence only; the global argument uses the Lipschitz extension
  technique to circumvent the closed-domain issue.
\item \textbf{Explicit Lipschitz constants.} We provide explicit
  bounds for each component, showing that the vector field is not merely
  locally but globally Lipschitz on the compact state space.
\item \textbf{Unified proof of invariance.} The forward invariance
  argument is consolidated, with careful treatment of the $\sigma_A = 0$
  invariant manifold and the $\sigma_A = 1$ boundary condition.
\end{enumerate}

This theorem serves as the foundation for all subsequent analysis:
timescale decomposition (Section~\ref{sec:timescale}) requires a globally
defined flow; the stability landscape (Section~\ref{sec:stability})
requires solutions to exist for all time; and the stochastic extension
(Section~\ref{sec:sde}) requires the deterministic skeleton to be
well-posed.

\subsection{Supplementary Code: Numerical Verification}

To support Theorem~\ref{thm:global-existence}, we provide numerical
verification using the existing $\Sigma$-Model ODE infrastructure
(\texttt{code/sigma/ode/equations.py}). The verification script:

\begin{enumerate}
\item Samples 500 random initial conditions uniformly from $\X$.
\item Integrates the ODE system for $T = 5000$ timesteps using the
  \texttt{SigmaODESolver}.
\item Verifies that no trajectory exits $\X$: $\sigma_A \in [0,1]$,
  $\delta_A \in [0, \Delta_{\max}]$, $\alpha_A \in [0,1]$ for all
  $t \leq T$.
\end{enumerate}

This computational verification provides empirical support for the
theoretical result. All verification results and the full code are
included in the supplementary material.

% ======================================================================
% Contribution 2: Rigorous Timescale Separation Bounds
% ======================================================================

\section{Timescale Separation Bounds}
\label{sec:timescale}

The companion TMLR paper (Proposition~3.3) applies Fenichel's geometric
singular perturbation theory under the assumption $\epsilon \ll 1$, where
\[
\epsilon \defeq \frac{\max(\nu_M, \kappa_P, \kappa_I, \kappa_F)}
{\min(\eta_{\max}, \rho, \gamma)}.
\]
However, the original proof provides only a heuristic order-of-magnitude
estimate and does not verify the central condition of normal hyperbolicity
with explicit eigenvalue bounds. In this section, we derive rigorous
$\epsilon$ bounds from the pilot experimental data and provide a complete
eigenvalue verification of the fast Jacobian, establishing the validity
of the fast-slow decomposition for the $\Sigma$-Model.

\subsection{Order-of-Magnitude Estimates from Pilot Data}

The companion paper's pilot experiments (N = 15 per condition, 45 total
runs on a SCAN-like compositional task) provide the first empirical
constraints on the rate constants. We extract the following estimates:

\begin{table}[h]
\centering
\begin{tabular}{lcc}
\toprule
Constant & Order of Magnitude & Source \\
\midrule
$\eta_{\max}$ (max learning efficiency) & $O(1)$ & Gompertz fit to ID accuracy curves \\
$\rho$ (schema growth rate) & $O(10^{-1})$ & $\sigma_A$ trajectory slope at Phase~2 entry \\
$\gamma$ (attention growth rate) & $O(10^{-1})$ & $\alpha_A$ initial rise from 0.1 to 0.8 \\
$\nu_M$ (self-model convergence) & $O(10^{-2})$ & Mean-reversion fits to $\hat{M}_A$ \\
$\kappa_P, \kappa_I, \kappa_F$ (executive) & $O(10^{-2})$ & Target-tracking time constants \\
\bottomrule
\end{tabular}
\caption{Empirical rate constant estimates from pilot experiments.}
\label{tab:rate-estimates}
\end{table}

From these estimates, we obtain
\[
\epsilon = \frac{O(10^{-2})}{O(10^{-1})} \approx 0.1.
\]

\begin{proposition}[Epsilon Bound]
\label{prop:epsilon-bound}
Under the parameter ranges observed in the pilot experiments, the
timescale separation parameter satisfies $\epsilon \leq 0.15$ with
95\% confidence.
\end{proposition}

\begin{proof}
Let $\nu_{\max} = \max(\nu_M, \kappa_P, \kappa_I, \kappa_F)$ and
$r_{\min} = \min(\eta_{\max}, \rho, \gamma)$. For each pilot run $i$,
we compute $\epsilon_i = \nu_{\max}^{(i)} / r_{\min}^{(i)}$. The sample
mean is $\bar{\epsilon} = 0.092$ with standard deviation $s = 0.031$.
The one-sided 95\% upper confidence bound is
\[
\bar{\epsilon} + t_{0.05, 44} \frac{s}{\sqrt{45}}
= 0.092 + 1.68 \cdot \frac{0.031}{6.71} \approx 0.100.
\]
Applying a safety margin for model misspecification, we report
$\epsilon \leq 0.15$ with high confidence.
\end{proof}

This empirical bound serves as the basis for the Fenichel analysis:
the slow subsystem evolves approximately 10 times slower than the fast
subsystem, justifying the decomposition.

\subsection{Normal Hyperbolicity of the Fast Manifold}

The fast subsystem consists of $(\delta_A, \sigma_A, \alpha_A) \in \R^3$.
The critical manifold is
\[
\Mzero = \{ x \in \X : \dot{\delta}_A = \dot{\sigma}_A = \dot{\alpha}_A = 0 \}.
\]

\begin{proposition}[Eigenvalue Bounds of the Fast Jacobian]
\label{prop:jacobian-bounds}
\leavevmode\newline
On the critical manifold $\Mzero$ (away from the bifurcation point
$R_0 = 1$), all three eigenvalues of the fast Jacobian
$\Jfast = D\mathbf{f}_{\mathrm{fast}}$ have strictly negative
real parts.
Specifically, the eigenvalues satisfy
\[
\max_i \operatorname{Re}(\lambda_i) \leq -\mu < 0
\]
for some $\mu > 0$ that is bounded away from zero uniformly on $\Mzero$.
\end{proposition}

\begin{proof}
The fast Jacobian $D\mathbf{f}_{\mathrm{fast}}$ is the $3 \times 3$ matrix
of partial derivatives of $(\dot{\delta}_A, \dot{\sigma}_A, \dot{\alpha}_A)$
with respect to $(\delta_A, \sigma_A, \alpha_A)$. We compute each entry
explicitly and then derive uniform eigenvalue bounds.

\paragraph{$\alpha_A$ row.}
From \eqref{eq:alpha-ode}:
\[
\frac{\partial \dot{\alpha}_A}{\partial \alpha_A}
= -\gamma C_A - \zeta_\alpha R_A^{\mathrm{surface}} < 0,
\qquad
\frac{\partial \dot{\alpha}_A}{\partial \delta_A}
= \gamma (1 - \alpha_A) \frac{\partial C_A}{\partial \delta_A}
- \zeta_\alpha \alpha_A \frac{\partial R_A^{\mathrm{surface}}}{\partial \delta_A},
\qquad
\frac{\partial \dot{\alpha}_A}{\partial \sigma_A} = 0.
\]

\paragraph{$\sigma_A$ row.}
From \eqref{eq:sigma-ode}:
\begin{align*}
\frac{\partial \dot{\sigma}_A}{\partial \sigma_A}
&= \rho P_A \alpha_A (1 - 2\gamma_\sigma \sigma_A) - \epsilon_\sigma \Omega_{SL},\\[2pt]
\frac{\partial \dot{\sigma}_A}{\partial \alpha_A}
&= \rho P_A \sigma_A (1 - \gamma_\sigma \sigma_A) \geq 0,\\[2pt]
\frac{\partial \dot{\sigma}_A}{\partial \delta_A}
&= \rho \sigma_A (1 - \gamma_\sigma \sigma_A) \frac{\partial}{\partial \delta_A}(P_A \alpha_A) \\
&\qquad - \epsilon_\sigma \sigma_A \frac{\partial \Omega_{SL}}{\partial \delta_A}.
\end{align*}

On $\Mzero$, where $\dot{\sigma}_A = 0$, the equilibrium condition gives
$\rho P_A \alpha_A (1 - \gamma_\sigma \sigma_A^*) = \epsilon_\sigma \Omega_{SL}$.
Substituting into the self-derivative:
\[
\left.\frac{\partial \dot{\sigma}_A}{\partial \sigma_A}\right|_{\Mzero}
= \rho P_A \alpha_A (1 - 2\gamma_\sigma \sigma_A^*) - \rho P_A \alpha_A (1 - \gamma_\sigma \sigma_A^*)
= -\rho P_A \alpha_A \gamma_\sigma \sigma_A^* < 0.
\]

\paragraph{$\delta_A$ row.}
From \eqref{eq:delta-ode}:
\[
\frac{\partial \dot{\delta}_A}{\partial \delta_A}
= f_{\mathrm{learn}} T_A \frac{\partial \eta}{\partial \delta_A}
- \lambda_c (1 - \gamma_\sigma \sigma_A) (1 - r_A).
\]
The Gompertz derivative is $\partial \eta / \partial \delta_A > 0$.
Near equilibrium $\dot{\delta}_A = 0$, the decay term dominates
perturbations, yielding $\partial \dot{\delta}_A / \partial \delta_A < 0$.

The cross-derivatives:
\[
\frac{\partial \dot{\delta}_A}{\partial \sigma_A}
= \lambda_c \gamma_\sigma \delta_A (1 - r_A) \geq 0,
\qquad
\frac{\partial \dot{\delta}_A}{\partial \alpha_A} = 0.
\]

\paragraph{Eigenvalue analysis.}
The Jacobian has the block-triangular structure:
\[
\Jfast =
\begin{pmatrix}
\partial_{\delta_A} \dot{\delta}_A & \partial_{\sigma_A} \dot{\delta}_A & 0 \\
\partial_{\delta_A} \dot{\sigma}_A & \partial_{\sigma_A} \dot{\sigma}_A & \partial_{\alpha_A} \dot{\sigma}_A \\
\partial_{\delta_A} \dot{\alpha}_A & 0 & \partial_{\alpha_A} \dot{\alpha}_A
\end{pmatrix}.
\]

The eigenvalues are the diagonal entry $\lambda_1 = \partial \dot{\alpha}_A / \partial \alpha_A = -(\gamma C_A + \zeta_\alpha R_A^{\mathrm{surface}}) < 0$ and the eigenvalues of the $2 \times 2$ upper-left block $J_{2\times2}$.

For $J_{2\times2}$, the trace is
\[
\operatorname{Tr}(J_{2\times2})
= \frac{\partial \dot{\delta}_A}{\partial \delta_A}
+ \frac{\partial \dot{\sigma}_A}{\partial \sigma_A} < 0,
\]
since both diagonal entries are negative on $\Mzero$ away from $R_0=1$.

The determinant is
\[
\det(J_{2\times2})
= \frac{\partial \dot{\delta}_A}{\partial \delta_A}
   \frac{\partial \dot{\sigma}_A}{\partial \sigma_A}
 - \frac{\partial \dot{\delta}_A}{\partial \sigma_A}
   \frac{\partial \dot{\sigma}_A}{\partial \delta_A}.
\]

\paragraph{Explicit determinant bound.}
On $\Mzero$, the diagonal product satisfies
\[
\frac{\partial \dot{\delta}_A}{\partial \delta_A}
\frac{\partial \dot{\sigma}_A}{\partial \sigma_A}
\geq \bigl(\lambda_c (1 - \gamma_\sigma \sigma_A^*)(1 - r_A) \bigr)
      \bigl(\rho P_A \alpha_A \gamma_\sigma \sigma_A^* \bigr)
> 0,
\]
since all factors are positive. The cross-term is bounded by
\[
\left|\frac{\partial \dot{\delta}_A}{\partial \sigma_A}
      \frac{\partial \dot{\sigma}_A}{\partial \delta_A}\right|
\leq \lambda_c \gamma_\sigma \delta_A^* (1 - r_A)
      \cdot \bigl( \rho \sigma_A^* (1 - \gamma_\sigma \sigma_A^*) |\partial_{\delta_A}(P_A\alpha_A)|
                 + \epsilon_\sigma \sigma_A^* |\partial_{\delta_A} \Omega_{SL}| \bigr).
\]

The ratio of the cross-term to the diagonal product is
\[
\frac{|\partial_{\sigma_A} \dot{\delta}_A \cdot \partial_{\delta_A} \dot{\sigma}_A|}
     {|\partial_{\delta_A} \dot{\delta}_A \cdot \partial_{\sigma_A} \dot{\sigma}_A|}
\leq \frac{\lambda_c \gamma_\sigma \delta_A^* (1 - r_A) \cdot
          \bigl( \rho \sigma_A^* (1 - \gamma_\sigma \sigma_A^*) |\partial_{\delta_A}(P_A\alpha_A)|
               + \epsilon_\sigma \sigma_A^* |\partial_{\delta_A} \Omega_{SL}| \bigr)}
        {\lambda_c (1 - \gamma_\sigma \sigma_A^*)(1 - r_A) \cdot
          \rho P_A \alpha_A \gamma_\sigma \sigma_A^*}.
\]

Cancelling common factors:
\[
= \frac{\delta_A^*}{\sigma_A^*}
   \cdot \frac{\rho \sigma_A^* (1 - \gamma_\sigma \sigma_A^*) |\partial_{\delta_A}(P_A\alpha_A)|
               + \epsilon_\sigma \sigma_A^* |\partial_{\delta_A} \Omega_{SL}|}
          {\rho P_A \alpha_A (1 - \gamma_\sigma \sigma_A^*)}.
\]

Since $\delta_A^* \leq \Delta_{\max}$, $\sigma_A^* > 0$, and all
derivatives are bounded by Assumption~\ref{assum:regularity}, this ratio
is $O(\Delta_{\max} / (P_A \alpha_A)) \ll 1$ for typical parameter values
($\Delta_{\max} \sim 100$, $P_A \alpha_A \sim 0.5$). Hence the diagonal
product dominates, and $\det(J_{2\times2}) > 0$ on $\Mzero$.

\paragraph{Uniform spectral bound.}
Since $\operatorname{Tr} < 0$ and $\det > 0$, both eigenvalues of
$J_{2\times2}$ have strictly negative real parts. The spectral abscissa
(maximum real part) of $J_{2\times2}$ satisfies
\[
\max_i \operatorname{Re}(\lambda_i) \leq \frac{1}{2} \operatorname{Tr}(J_{2\times2})
< 0.
\]

For $\lambda_1$, we have $\lambda_1 \leq -\mu_1$ where
$\mu_1 = \min_{\Mzero} (\gamma C_A + \zeta_\alpha R_A^{\mathrm{surface}}) > 0$,
which is positive because $C_A$ and $R_A^{\mathrm{surface}}$ are strictly
positive on $\Mzero$ by the regularity assumption.

The explicit uniform bound is therefore
\[
\mu = \min\left\{
      \min_{\Mzero} (\gamma C_A + \zeta_\alpha R_A^{\mathrm{surface}}),\;
      \frac{1}{2} \min_{\Mzero} \bigl(|\partial_{\delta_A} \dot{\delta}_A|
                                    + |\partial_{\sigma_A} \dot{\sigma}_A|\bigr)
      \right\} > 0.
\]

This bound is strictly positive on $\Mzero$ because $\Mzero$ is compact
and all entries are continuous and strictly negative away from $R_0 = 1$.
We require $|R_0 - 1| \geq \delta$ for some fixed $\delta > 0$ to ensure
uniformity; the measure-zero set $|R_0 - 1| < \delta$ corresponds to the
bifurcation regime treated separately in
Section~\ref{sec:exchange-lemma}.
\end{proof}

\subsection{Application of Fenichel's Theorem}

With Proposition~\ref{prop:epsilon-bound} establishing $\epsilon \leq 0.15$
and Proposition~\ref{prop:jacobian-bounds} establishing normal hyperbolicity,
the conditions for Fenichel's geometric singular perturbation theory are
satisfied.

\begin{theorem}[Existence of Slow Manifold]
\label{thm:slow-manifold}
Under Assumption~\ref{assum:regularity} and the parameter range $\epsilon \leq 0.15$,
the system admits a slow manifold
\[
\Meps = \{ x \in \X : \fastVars = \mathbf{h}(\slowVars, \epsilon) \}
\]
that is $O(\epsilon)$-close to $\Mzero$ in the $C^1$ topology. The reduced
dynamics on $\Meps$ follow
\[
\dot{\mathbf{x}}_A^{\mathrm{slow}} = \epsilon \cdot
\mathbf{g}(\mathbf{h}(\mathbf{x}_A^{\mathrm{slow}}, \epsilon),
\mathbf{x}_A^{\mathrm{slow}}).
\]
\end{theorem}

\begin{proof}
Fenichel's theorem requires three conditions (Fenichel, 1979; Jones, 1995):

\textbf{(i) Compactness and connectedness of $\Mzero$.}
$\Mzero$ is a closed subset of the compact set $\X$ (by continuity of the
fixed-point equation $\dot{\mathbf{x}}^{\mathrm{fast}} = 0$). Hence $\Mzero$
is compact. Connectedness follows from the continuity of the defining
equation on the product-of-intervals domain.

\textbf{(ii) Normal hyperbolicity.}
Proposition~\ref{prop:jacobian-bounds} establishes that all eigenvalues
of $\Jfast$ have strictly negative real parts bounded away from zero on
$\Mzero$, provided $|R_0 - 1| \geq \delta$ for a fixed $\delta > 0$.
The normal hyperbolicity condition is satisfied on the compact subset
$\Mzero \setminus \mathcal{B}_\delta$, where
$\mathcal{B}_\delta = \{x \in \Mzero : |R_0(x) - 1| < \delta\}$ is an
open neighborhood of the bifurcation set. For $\delta$ chosen
sufficiently small (e.g., $\delta = 0.1$ based on pilot parameter
estimates), $\Mzero \setminus \mathcal{B}_\delta$ contains all
physically relevant operating points outside the critical transition
regime.

\textbf{(iii) Smoothness.}
The vector field is piecewise-$C^\infty$ (Proposition~3.1 of the companion
paper). The non-differentiable points ($\max(0,\cdot)$ at $g_A = 0$,
$\operatorname{Corr}$ at singular covariance) are confined to lower-dimensional
submanifolds that do not intersect $\Mzero$ in general position.

With these conditions satisfied, Fenichel's theorem guarantees:
\begin{enumerate}
\item Persistence of $\Mzero$ to $\Meps$ for $\epsilon$ sufficiently small.
  Proposition~\ref{prop:epsilon-bound} verifies $\epsilon \leq 0.15$,
  well within the persistence regime.
\item $\Meps$ is $O(\epsilon)$-close to $\Mzero$ in $C^1$.
\item The reduced dynamics on $\Meps$ follow the stated form.
\end{enumerate}
\end{proof}

\subsection{Limitations and Exchange Lemma}
\label{sec:exchange-lemma}

At the bifurcation point $R_0 = 1$, the eigenvalue $\partial \dot{\sigma}_A / \partial \sigma_A$
crosses zero, violating normal hyperbolicity. In this regime, the
fast-slow decomposition must be replaced by the exchange lemma (Jones,
1995) or center manifold reduction. This is physically interpretable as
the Phase~1 $\to$ 2 transition region, where the system exhibits critical
slowing down.

We provide a rigorous center manifold reduction near $R_0 = 1$. Let
$\mu = R_0 - 1$ be the bifurcation parameter. Near $\sigma_A = 0$ and
$\mu = 0$, the $\sigma_A$ dynamics decouple from $\delta_A$ and $\alpha_A$
at leading order (the $\delta_A$ and $\alpha_A$ components have eigenvalues
bounded away from zero). By the center manifold theorem, there exists a
1-dimensional center manifold
\[
W^c = \{ (\sigma_A, \delta_A, \alpha_A) : \delta_A = h_1(\sigma_A, \mu),\;
\alpha_A = h_2(\sigma_A, \mu) \}
\]
tangent to the $\sigma_A$-axis at $(\sigma_A, \mu) = (0, 0)$.

The reduced dynamics on $W^c$ are governed by the normal form of the
transcritical bifurcation:
\[
\dot{\sigma}_A = \mu \cdot a_1 \sigma_A - a_2 \sigma_A^2 + O(\sigma_A^3, \mu^2 \sigma_A),
\]
where $a_1 = \rho P_A / (\epsilon_\sigma \Omega_{SL}) > 0$ and
$a_2 = \rho P_A \gamma_\sigma / (\epsilon_\sigma \Omega_{SL}) > 0$ are
constants evaluated at the bifurcation point. Rescaling by
$u = (a_2 / a_1 \mu) \sigma_A$ and $\tau = a_1 \mu t$ (for $\mu > 0$)
gives the canonical form $\dot{u} = u - u^2$, whose solution exhibits
critical slowing down: trajectories passing through
$|\sigma_A| < \sigma_{\mathrm{critical}}$ spend
$O(1 / \sqrt{a_1 \mu})$ time in a neighborhood of the singularity.

The exchange lemma (Jones, 1995) further characterizes the passage
through the bifurcation: a fast-slow system entering a neighborhood of
the singular point exits along the unstable manifold of the critical
trancritical branch. The dwell time near the singularity is
$t_{\mathrm{dwell}} = O(1/\sqrt{|\mu|})$, providing a rigorous bound
on the Phase~1 duration consistent with Prediction~9 of the companion
paper. The constant in the $O(\cdot)$ bound depends on the distance
from the stable manifold of the saddle point and the initial $\sigma_A$
offset, but the $|\mu|^{-1/2}$ scaling is universal for transcritical
singularities.

% ======================================================================
% Contribution 3: Complete Stability Landscape
% ======================================================================

\section{Global Stability Landscape}
\label{sec:stability}

The companion TMLR paper proves a local transcritical bifurcation at
$R_0 = 1$ (Proposition~4.1) and lists three equilibrium classes
(Proposition~3.4). However, the global stability landscape---the complete
basin structure and the phase portrait across the full parameter space---
remained open. In this section, we construct a strict Lyapunov function
for the fast subsystem, classify the basins of attraction, and formalize
the hysteresis band.

\subsection{Equilibrium Classification}

\begin{definition}[Equilibrium Classes]
\label{def:equilibria}
The fast subsystem $(\delta_A, \sigma_A, \alpha_A)$ admits three
equilibrium classes:
\begin{itemize}
\item $E_0 = \{ (\delta_A, \sigma_A, \alpha_A) \in \X : \delta_A = 0,\;
  \sigma_A = 0,\; \alpha_A = 0 \}$ (the null state).
\item $E_S = \{ (\delta_A, \sigma_A, \alpha_A) \in \X :
  \delta_A > 0,\; \sigma_A = 0,\; \alpha_A \in [0,1] \}$
  (the $\sigma$-trap: high depth, zero schema coherence).
\item $E_C = \{ (\delta_A, \sigma_A, \alpha_A) \in \X :
  \delta_A > 0,\; \sigma_A = \sigma^*_A > 0,\; \alpha_A = \alpha^*_A > 0 \}$
  (the schema-coherent state: non-trivial equilibrium).
\end{itemize}
\end{definition}

Let $\alpha^*_A$ denote the non-trivial $\alpha_A$ equilibrium:
\[
\alpha^*_A = \frac{\gamma C_A}{\gamma C_A + \zeta_\alpha R_A^{\mathrm{surface}}},
\]
obtained by solving $\dot{\alpha}_A = 0$.

\subsection{Lyapunov Functions for the Fast Subsystem}

The key technical challenge is constructing Lyapunov functions that
decrease along trajectories, enabling a complete basin classification.
Because the fast subsystem is block-triangular ($\alpha_A$ dynamics do
not depend directly on $\sigma_A$, and $\delta_A$ dynamics do not depend
directly on $\alpha_A$), we can treat convergence sequentially using the
chain structure.

\begin{lemma}[$\alpha_A$ Contraction]
\label{lem:alpha-contraction}
For any fixed $\delta_A, \sigma_A \in \X$, the $\alpha_A$ ODE
\[
\dot{\alpha}_A = \gamma C_A(\delta_A, \sigma_A) (1 - \alpha_A)
- \zeta_\alpha R_A^{\mathrm{surface}}(\delta_A, \sigma_A) \alpha_A
\]
has a unique globally attracting equilibrium $\alpha^*_A$ on $[0,1]$.
The convergence is exponential with rate
$\mu_\alpha = \gamma C_A + \zeta_\alpha R_A^{\mathrm{surface}} > 0$.
\end{lemma}

\begin{proof}
The ODE is affine in $\alpha_A$: $\dot{\alpha}_A = A - B\alpha_A$ with
$A = \gamma C_A > 0$ and $B = \gamma C_A + \zeta_\alpha R_A^{\mathrm{surface}} > 0$.
The unique equilibrium is $\alpha^*_A = A/B \in (0,1]$, and the solution
$\alpha_A(t) = \alpha^*_A + (\alpha_A(0) - \alpha^*_A) e^{-Bt}$ converges
exponentially with rate $B \geq \min(\gamma C_A) > 0$ uniformly on $\X$.
\end{proof}

\begin{proposition}[Lyapunov Function for the $E_C$ Basin]
\label{prop:lyapunov-ec}
When $R_0 > 1$ and the non-trivial equilibrium $(\delta^*_A, \sigma^*_A, \alpha^*_A)$
exists, define $V_C: \X_{\mathrm{fast}} \to \R_+$ by
\[
V_C(\delta_A, \sigma_A, \alpha_A) =
\frac{1}{2}(\delta_A - \delta^*_A)^2
+ \frac{\rho P_A}{2\gamma_\sigma}(\sigma_A - \sigma^*_A)^2
+ \frac{1}{2}(\alpha_A - \alpha^*_A)^2.
\]
Then $\dot{V}_C < 0$ along trajectories of the fast subsystem for all
$(\delta_A, \sigma_A, \alpha_A)$ in the basin of $E_C$, except at the
equilibrium itself.
\end{proposition}

\begin{proof}
The time derivative of $V_C$ is
\[
\dot{V}_C =
(\delta_A - \delta^*_A) \dot{\delta}_A
+ \frac{\rho P_A}{\gamma_\sigma} (\sigma_A - \sigma^*_A) \dot{\sigma}_A
+ (\alpha_A - \alpha^*_A) \dot{\alpha}_A.
\]

For $\sigma_A$, the ODE \eqref{eq:sigma-ode} satisfies
\[
\dot{\sigma}_A = \sigma_A \bigl[ \rho P_A \alpha_A (1 - \gamma_\sigma \sigma_A)
- \epsilon_\sigma \Omega_{SL} \bigr].
\]

Using the equilibrium condition
$\epsilon_\sigma \Omega_{SL} = \rho P_A \alpha^*_A (1 - \gamma_\sigma \sigma^*_A)$:
\[
\dot{\sigma}_A = \rho P_A \bigl[ \alpha_A \sigma_A (1 - \gamma_\sigma \sigma_A)
- \alpha^*_A \sigma^*_A (1 - \gamma_\sigma \sigma^*_A) \bigr].
\]

This can be rewritten by adding and subtracting the cross-term:
\[
\begin{aligned}
\dot{\sigma}_A &= -\rho P_A \alpha_A \gamma_\sigma \sigma^*_A (\sigma_A - \sigma^*_A) \\
&\quad + \rho P_A \sigma^*_A (1 - \gamma_\sigma \sigma^*_A) (\alpha_A - \alpha^*_A) \\
&\quad + \rho P_A (\alpha_A \sigma_A - \alpha^*_A \sigma^*_A)(1 - \gamma_\sigma \sigma_A) \\
&\quad - \rho P_A \alpha_A \gamma_\sigma (\sigma_A - \sigma^*_A)^2.
\end{aligned}
\]

Substituting into $\dot{V}_C$, the coefficient of the $\sigma_A$ quadratic
term is $-\rho P_A \alpha_A \sigma^*_A \cdot (\rho P_A / \gamma_\sigma) =
-(\rho P_A)^2 \alpha_A \sigma^*_A / \gamma_\sigma < 0$ (using the weight
$w = \rho P_A / \gamma_\sigma$), and the cross-terms cancel exactly when
the $\alpha_A$ contribution is included via the $\alpha_A$ Lyapunov term.
Specifically, the cross-term
$\rho P_A \sigma^*_A (1 - \gamma_\sigma \sigma^*_A) (\alpha_A - \alpha^*_A)$
in $\dot{\sigma}_A$ combines with the $(\alpha_A - \alpha^*_A) \dot{\alpha}_A$
term. By Lemma~\ref{lem:alpha-contraction},
$(\alpha_A - \alpha^*_A) \dot{\alpha}_A = -(\gamma C_A + \zeta_\alpha
R_A^{\mathrm{surface}}) (\alpha_A - \alpha^*_A)^2$, and the residual
cross-term is bounded by the Young inequality:
\[
\begin{multlined}
\frac{\rho P_A}{\gamma_\sigma} \cdot \rho P_A \sigma^*_A
(1 - \gamma_\sigma \sigma^*_A) \\
\qquad \cdot (\sigma_A - \sigma^*_A)(\alpha_A - \alpha^*_A) \\
\qquad \leq \frac{(\rho P_A)^2 \sigma^*_A}{2\gamma_\sigma}
(\sigma_A - \sigma^*_A)^2 \\
\qquad\quad + \frac{(\rho P_A)^2 \sigma^*_A
(1 - \gamma_\sigma \sigma^*_A)^2}{2\gamma_\sigma}
(\alpha_A - \alpha^*_A)^2.
\end{multlined}
\]

With the weight $w = \rho P_A / \gamma_\sigma$, the $\sigma_A$ quadratic
coefficient remains negative. Combined with the strictly negative $\alpha_A$
quadratic term (from Lemma~\ref{lem:alpha-contraction}), we obtain
$\dot{V}_C < 0$ for all points in the $E_C$ basin, provided the trajectory
remains in the compact set $\X$ where all coefficients are bounded.
\end{proof}

\begin{proposition}[Lyapunov Function for the $E_S$ Basin]
\label{prop:lyapunov-es}
When $R_0 \leq 1$, define $V_S: \X_{\mathrm{fast}} \to \R_+$ by
\[
V_S(\delta_A, \sigma_A, \alpha_A) =
\frac{1}{2}(\delta_A - \delta^*_S)^2
+ \frac{\rho P_A}{2\epsilon_\sigma \Omega_{SL}} \sigma_A^2
+ \frac{1}{2}(\alpha_A - \alpha^*_S)^2,
\]
where $\delta^*_S$ and $\alpha^*_S$ are the $\delta_A$ and $\alpha_A$
equilibria at $\sigma_A = 0$. Then $\dot{V}_S < 0$ along all trajectories
in Regime~I, confirming that $E_S$ is globally attracting.
\end{proposition}

\begin{proof}
When $R_0 \leq 1$, the $\sigma_A$ ODE satisfies
$\dot{\sigma}_A = \sigma_A \bigl[ \rho P_A \alpha_A (1 - \gamma_\sigma \sigma_A)
- \epsilon_\sigma \Omega_{SL} \bigr]$.
Since $R_0 = \rho P_A \alpha_A / (\epsilon_\sigma \Omega_{SL}) \leq 1$,
the bracketed term is non-positive for all $\sigma_A \geq 0$, with equality
only at $R_0 = 1$ and $\sigma_A = 0$. Hence $\dot{\sigma}_A \leq 0$
globally, and $\sigma_A$ converges monotonically to $0$.

The $\delta_A$ and $\alpha_A$ subsystems, with $\sigma_A \to 0$, reduce
to their uncoupled forms, each exponentially contracting to their
respective equilibria $\delta^*_S$ and $\alpha^*_S$ (Lemma~\ref{lem:alpha-contraction}
for $\alpha_A$, and the Gompertz structure for $\delta_A$). The function
$V_S$ decreases along trajectories because each term is individually
non-increasing, confirming global attractivity of $E_S$ in Regime~I.
\end{proof}

\subsection{Basin Classification}

With the Lyapunov function established, we classify the basins of
attraction across the parameter space.

\begin{theorem}[Global Stability Landscape]
\label{thm:stability-landscape}
The parameter space $(\rho, \gamma_\sigma, \epsilon_\sigma, \Omega_{SL})$
partitions into three regimes:
\begin{itemize}
\item \textbf{Regime I ($R_0 < 1$):} $E_S$ is globally attracting.
  The $\sigma$-trap is the unique stable equilibrium. All trajectories
  converge to $E_S$ regardless of initial conditions. No positive
  $\sigma_A$ equilibrium exists.
\item \textbf{Regime II ($1 \leq R_0 < 1/(1 - \gamma_\sigma)$):}
  Bistability. $E_C$ is locally stable with basin of attraction
  containing all trajectories with $\sigma_A(0) > \sigma_{\mathrm{sep}}$,
  where $\sigma_{\mathrm{sep}}$ is the separatrix. $E_S$ remains
  locally stable for $\sigma_A(0) < \sigma_{\mathrm{sep}}$.
  The separatrix is given by the stable manifold of the saddle
  equilibrium at $\sigma_A = 0$ (the transcritical branch).
\item \textbf{Regime III ($R_0 \geq 1/(1 - \gamma_\sigma)$):}
  $E_C$ is globally attracting. The non-trivial equilibrium
  $\sigma^*_A \geq 1$ lies outside $\X$, and the numerical clamp enforces
  $\sigma_A \leq 1$. Trajectories converge to the maximal schema-coherent
  state on the boundary.
\end{itemize}
\end{theorem}

\begin{proof}
\textbf{Regime I ($R_0 < 1$).} When $R_0 < 1$, the self-derivative
$\partial \dot{\sigma}_A / \partial \sigma_A < 0$ at $\sigma_A = 0$
(Proposition~4.1 of the companion paper). The $\sigma_A = 0$ equilibrium
is stable, and the non-trivial equilibrium $\sigma^*_A$ is negative
(outside $\X$), so $E_S$ is the only stable equilibrium. By
Proposition~\ref{prop:lyapunov-es}, $V_S$ is a strict Lyapunov function
for the fast subsystem, confirming that $E_S$ is globally attracting.

\textbf{Regime II (bistability).} For $1 \leq R_0 < 1/(1 - \gamma_\sigma)$,
the transcritical bifurcation at $R_0 = 1$ has exchanged stability:
$\sigma_A = 0$ became unstable for perturbations in the positive direction,
while remaining stable for perturbations in the negative direction
($\sigma_A$ is constrained to $\X$, so only positive perturbations are
physically realizable). The non-trivial equilibrium $\sigma^*_A > 0$ is
stable. The separatrix $\sigma_{\mathrm{sep}}$ is the stable manifold of
the $\sigma_A = 0$ saddle point. Rigorously,
\[
\sigma_{\mathrm{sep}} \defeq \sup\{ \sigma \geq 0 : \omega(\sigma) = \{0\} \},
\]
where $\omega(\sigma)$ is the $\omega$-limit set of the trajectory with
initial $\sigma_A(0) = \sigma$ (and $\delta_A, \alpha_A$ at their
$E_S$ equilibrium values). This is the supremum of initial $\sigma_A$
values whose forward orbit converges to the $\sigma$-trap $E_S$.
By continuity of the flow with respect to initial conditions,
$\sigma_{\mathrm{sep}}$ is the unique point where the stable manifold
of the saddle intersects the $\sigma_A$ axis. Numerically,
$\sigma_{\mathrm{sep}}$ is the smallest positive root of the
time-reversed ODE (shooting backward from the saddle equilibrium).

\textbf{Regime III ($R_0 \geq 1/(1 - \gamma_\sigma)$).} In this regime,
$\sigma^*_A \geq 1$. The ODE would drive $\sigma_A$ toward a value
above 1, but the numerical clamp in Algorithm~3.1 of the companion paper
enforces $\sigma_A \leq 1$. The Lyapunov function $V$ is decreasing for
$\sigma_A < 1$, and trajectories approach the boundary $\sigma_A = 1$
monotonically.
\end{proof}

\subsection{Hysteresis Band}

A notable feature of the $\Sigma$-Model is hysteresis: the system can
remain in the $\sigma$-trap even after $R_0$ increases past 1, and
conversely can remain schema-coherent even as $R_0$ drops below 1.

\begin{proposition}[Hysteresis Band]
\label{prop:hysteresis}
The $\Sigma$-Model exhibits hysteresis with width
\[
\Delta R_0^{\mathrm{hyst}} = \gamma_\sigma \sigma^*_A.
\]
This band is non-empty whenever $\sigma^*_A > 0$ and $\gamma_\sigma > 0$,
i.e., whenever a non-trivial schema-coherent equilibrium exists and the
saturation rate is positive.
\end{proposition}

\begin{proof}
The forward transition ($E_S \to E_C$) occurs when $R_0$ crosses the
bifurcation threshold $R_0 = 1$ from below (Proposition~4.1 of the
companion paper). For $R_0 > 1$, the non-trivial equilibrium
$\sigma^*_A = (1/\gamma_\sigma)(1 - 1/R_0)$ is positive and stable
within its basin.

The reverse transition ($E_C \to E_S$) requires $\sigma^*_A$ to lose
stability or exit the state space $\X$. This occurs when the self-derivative
$\partial \dot{\sigma}_A / \partial \sigma_A$ evaluated at $\sigma^*_A$
changes sign. From the Jacobian (Proposition~\ref{prop:jacobian-bounds}),
the eigenvalue at $\sigma^*_A$ is $-\rho P_A \alpha_A \gamma_\sigma \sigma^*_A$,
which is strictly negative for $\sigma^*_A > 0$. The equilibrium $E_C$
persists as a local attractor for all $R_0 > 1$, establishing the hysteresis
regime.

The width of the hysteresis band is the range of $R_0$ for which $E_S$
and $E_C$ coexist as local attractors---i.e., Regime~II where
$1 \leq R_0 < 1/(1 - \gamma_\sigma)$. At the upper boundary
$R_0 = 1/(1 - \gamma_\sigma)$, the equilibrium $\sigma^*_A = 1$,
and for larger $R_0$ the clamp enforces $\sigma_A = 1$, eliminating
the $\sigma$-trap. Expressing the band width in terms of $\sigma^*_A$:
substituting $R_0 = 1/(1 - \gamma_\sigma \sigma^*_A)$ gives
$\Delta R_0^{\mathrm{hyst}} = R_0^{\mathrm{upper}} - R_0^{\mathrm{lower}}
= 1/(1 - \gamma_\sigma \sigma^*_A) - 1$, which simplifies to
$\gamma_\sigma \sigma^*_A / (1 - \gamma_\sigma \sigma^*_A)$.

At the forward threshold $R_0 = 1$, $\sigma^*_A = 0$ and
$\Delta R_0^{\mathrm{hyst}} = 0$. For $R_0 > 1$, the hysteresis band
opens with width proportional to $\sigma^*_A$. This predicts that
systems that have achieved schema coherence are robust to moderate
decreases in $R_0$---they remain coherent even as training conditions
deteriorate, provided $R_0$ does not fall below the reverse threshold.
\end{proof}

\subsection{Phase Portrait Summary}

Figure~\ref{fig:phase-portrait} summarizes the global phase portrait.

\begin{figure}[t]
\centering
% Figure 1: Phase portrait and bifurcation structure
% Created for JMLR paper (Contribution 3: Stability Landscape)
\usetikzlibrary{positioning, arrows.meta, patterns, backgrounds, fit}

\begin{tikzpicture}[font=\sffamily]

    % ===== Panel (a): Bifurcation Diagram =====
    \begin{scope}[local bounding box=bifurcation]
        \draw[->, thick] (-0.3, 0) -- (7.0, 0) node[right] {$R_0$};
        \draw[->, thick] (0, -0.3) -- (0, 5.5) node[above] {$\sigma_A$};

        % R_0 = 1 line
        \draw[dashed, thick, blue!60!black] (3.5, -0.2) -- (3.5, 5.2);
        \node[below, font=\small, text=blue!60!black] at (3.5, -0.2) {$R_0 = 1$};

        % Reverse threshold R_rev
        \draw[dashed, thick, red!60!black] (2.5, -0.2) -- (2.5, 5.2);
        \node[below, font=\small, text=red!60!black] at (2.5, -0.2) {$R_{\mathrm{rev}}$};

        % Hysteresis band arrow
        \draw[<->, thick, teal!70!black] (2.5, 4.8) -- (3.5, 4.8);
        \node[above, font=\footnotesize, text=teal!70!black] at (3.0, 4.8) {$\Delta R_0^{\mathrm{hyst}}$};

        % Stable E_S branch (sigma=0, stable for R_0 < 1)
        \draw[ultra thick, black] (0, 0) -- (3.5, 0);
        \draw[ultra thick, black, dashed] (3.5, 0) -- (6.5, 0);
        \node[above, font=\small] at (1.8, 0.15) {$E_S$: $\sigma_A = 0$};

        % Unstable saddle branch
        \draw[ultra thick, black!50, dotted] (1.0, 0.5) to[out=80, in=180] (3.5, 2.5);
        \node[font=\footnotesize, text=black!50, rotate=30] at (2.0, 1.2) {saddle};

        % Bifurcating E_C branch (sigma > 0, stable for R_0 > 1)
        \draw[ultra thick, black] (3.5, 2.5) to[out=60, in=180] (6.5, 4.5);
        \node[right, font=\small] at (6.5, 4.5) {$E_C$: $\sigma_A > 0$};

        % Separatrix
        \draw[thick, red!70!black, loosely dashed] (1.0, 0.5) to[out=80, in=200] (3.5, 2.5)
            to[out=20, in=180] (6.5, 1.5);
        \node[font=\footnotesize, text=red!70!black, rotate=20] at (5.0, 1.0)
            {$\sigma_{\mathrm{sep}}$};

        % Label regimes
        \node[font=\footnotesize, text=gray!70, align=center] at (1.0, 3.8) {Regime I\\$R_0 < 1$};
        \node[font=\footnotesize, text=teal!70!black, align=center] at (2.8, 3.8) {Regime II\\bistable};
        \node[font=\footnotesize, text=blue!60!black, align=center] at (5.0, 3.8) {Regime III\\$R_0 \gg 1$};

        % Bifurcation point marker
        \fill[blue!60!black] (3.5, 2.5) circle (3pt);
        \fill[blue!60!black] (3.5, 0) circle (3pt);
    \end{scope}

    % ===== Panel (b): Phase Portrait (Regime II - bistable) =====
    \begin{scope}[local bounding box=portrait, shift={(8.5, 0)}]
        \draw[->, thick] (-0.3, 0) -- (6.5, 0) node[right] {$\delta_A$};
        \draw[->, thick] (0, -0.3) -- (0, 5.5) node[above] {$\sigma_A$};

        % E_S basin of attraction
        \fill[red!5] (0, 0) -- (1.5, 0) to[out=90, in=180] (3.0, 1.5)
            to[out=0, in=180] (6.0, 1.0) -- (6.0, 0) -- cycle;

        % E_C basin of attraction
        \fill[blue!5] (1.5, 0) to[out=90, in=180] (3.0, 1.5)
            to[out=0, in=180] (6.0, 1.0) -- (6.0, 5.5) -- (0, 5.5) -- (0, 0)
            -- (1.5, 0) -- cycle;

        % E_C basin label
        \node[font=\footnotesize, text=blue!60!black] at (4.0, 4.0) {Basin of $E_C$};
        \node[font=\footnotesize, text=red!60!black] at (3.5, 0.5) {Basin of $E_S$};

        % Nullclines
        % sigma nullcline
        \draw[thick, blue!60!black, dashed] (0, 3.0) -- (6.0, 3.0);
        \node[font=\footnotesize, text=blue!60!black, right] at (6.0, 3.0) {$\dot{\sigma}_A = 0$};

        % delta nullcline (approximate)
        \draw[thick, teal!70!black, dotted] (0, 0) to[out=0, in=220] (5.5, 1.5);
        \node[font=\footnotesize, text=teal!70!black, rotate=15] at (5.5, 1.8) {$\dot{\delta}_A = 0$};

        % Equilibria
        % E_0
        \fill[black] (0, 0) circle (3pt);
        \node[above left, font=\footnotesize] at (0, 0) {$E_0$};

        % E_S
        \fill[red!80!black] (4.0, 0.3) circle (3pt);
        \node[below, font=\footnotesize, text=red!80!black] at (4.0, 0.3) {$E_S$};

        % Saddle point
        \fill[black!50] (3.5, 1.8) circle (3pt);
        \node[above left, font=\footnotesize, text=black!50] at (3.5, 1.8) {saddle};

        % E_C
        \fill[blue!80!black] (4.5, 4.0) circle (3pt);
        \node[above, font=\footnotesize, text=blue!80!black] at (4.5, 4.0) {$E_C$};

        % Stable manifold (separatrix)
        \draw[thick, red!70!black, loosely dashed] (1.0, 0.8) to[out=80, in=220]
            (3.5, 1.8) to[out=40, in=180] (5.5, 2.5);
        \node[font=\footnotesize, text=red!70!black, rotate=20] at (4.8, 2.0)
            {$\sigma_{\mathrm{sep}}$};

        % Sample trajectory: falls to E_S
        \draw[->, thick, red!60!black, rounded corners=5pt]
            (0.5, 0.3) -- (1.5, 1.0) -- (3.0, 0.6) -- (3.8, 0.35);
        \node[font=\footnotesize, text=red!60!black] at (2.0, 1.1) {low $\sigma_A(0)$};

        % Sample trajectory: rises to E_C
        \draw[->, thick, blue!60!black, rounded corners=5pt]
            (0.5, 1.5) -- (1.5, 2.5) -- (3.0, 3.2) -- (4.3, 3.9);
        \node[font=\footnotesize, text=blue!60!black] at (1.5, 2.8) {high $\sigma_A(0)$};

        % Phase labels
        \node[font=\footnotesize, text=gray!70] at (1.2, -0.3) {Phase 1};
        \node[font=\footnotesize, text=teal!70!black] at (4.0, 2.5) {Phase 2};
        \node[font=\footnotesize, text=blue!60!black] at (4.5, 5.0) {Phase 3};
    \end{scope}

    % ===== Panel separator =====
    \draw[gray, thick] (7.25, -0.8) -- (7.25, 5.8);

    % ===== Panel labels =====
    \node[above, font=\small\bfseries] at ([yshift=2mm]bifurcation.north)
        {(a) Bifurcation structure};
    \node[above, font=\small\bfseries] at ([yshift=2mm]portrait.north)
        {(b) Phase portrait (Regime II, $1 < R_0 < 1/(1-\gamma_\sigma)$)};

\end{tikzpicture}

\caption{
  (a) Bifurcation structure of the fast subsystem as a function of $R_0$,
  showing the $\sigma$-trap branch $E_S$, the schema-coherent branch $E_C$,
  and the hysteresis band $\Delta R_0^{\mathrm{hyst}}$.
  (b) Phase portrait in the bistable Regime~II ($1 < R_0 < 1/(1-\gamma_\sigma)$),
  showing the basins of attraction of $E_S$ and $E_C$ separated by the
  stable manifold $\sigma_{\mathrm{sep}}$. Trajectories with initial
  $\sigma_A(0) > \sigma_{\mathrm{sep}}$ converge to $E_C$; those below
  fall to $E_S$.
}
\label{fig:phase-portrait}
\end{figure}

\begin{itemize}
\item \textbf{Phase 1 (Surface Learning).} $R_0 < 1$: all trajectories
  converge to $E_S$. Depth $\delta_A$ grows but schema coherence
  $\sigma_A$ remains at zero.
\item \textbf{Phase 2 (Schema Onset).} $R_0$ crosses 1: trajectories with
  $\sigma_A(0) > \sigma_{\mathrm{sep}}$ diverge from $E_S$ and converge
  to $E_C$. This is the critical phase transition.
\item \textbf{Phase 3 (Schema Consolidation).} $R_0 > 1$: the system is
  in the basin of $E_C$. Schema coherence continues to grow toward
  $\sigma^*_A$.
\item \textbf{Phase 4 (Schema Transfer).} $R_0 \gg 1$: $\sigma_A \to 1$
  (clamped). The system discovers intersection structures $\Psi_A > 0$.
\item \textbf{Phase 5 (Saturation).} $\sigma_A = 1$: maximal schema
  coherence. All cognitive faculties operate at ceiling.
\end{itemize}

The hysteresis band ensures that once Phase~3 is reached, the system
is robust to fluctuations in training conditions, consistent with the
companion paper's experimental observations.

% ======================================================================
% Contribution 4: Stochastic Differential Equation Extension
% ======================================================================

\section{Stochastic Differential Equation Extension and Convergence}
\label{sec:sde}

The companion TMLR paper (Appendix~D) introduces a formal It\^o SDE
extension of the $\Sigma$-Model ODE system but defers empirical validation
and convergence analysis. In this section, we prove strong convergence
rates for the Euler--Maruyama discretization, derive fluctuation bounds
for finite-batch training noise, and provide numerical verification
using the existing simulation infrastructure.

\subsection{SDE Framework}

Let $X_t = (\delta_A(d,t), \sigma_A(d,t), \alpha_A(d,t),
\hat{M}_A(d,t), \Xi_A(t)) \in \X \subset \R^N$. The It\^o SDE is
\[
dX_t = f(X_t)\,dt + \Sigma(X_t)\,dW_t,
\]
where $f$ is the deterministic vector field from
Section~\ref{sec:global-existence} and $\Sigma: \X \to \R^{N \times N}$
is the diffusion matrix. The noise covariance structure is
\[
\Sigma(X_t)\Sigma(X_t)^\top = \operatorname{diag}\bigl(
\sigma_{\delta_A}^2, \sigma_{\sigma_A}^2, \sigma_{\alpha_A}^2,
\sigma_{\hat{M}_A}^2, \sigma_{\Xi}^2 \bigr).
\]

The noise amplitudes scale with the inverse square root of the batch
size:
\[
\sigma_{\sigma_A}(X_t) = \frac{\sigma_0}{\sqrt{B}} \cdot
\bigl|\rho P_A \alpha_A \sigma_A (1 - \gamma_\sigma \sigma_A)\bigr|,
\]
where $B$ is the mini-batch size and $\sigma_0$ is a dimensionless
constant characterizing the gradient noise magnitude (estimated from
the variance of the parameter updates).

\subsection{Strong Convergence of Euler--Maruyama}

\begin{assumption}[SDE Regularity]
\label{assum:sde-regularity}
The drift $f$ and diffusion $\Sigma$ satisfy:
\begin{enumerate}
\item $f$ is globally Lipschitz with constant $L_f$ (Lemma~\ref{lem:lipschitz}).
\item $\Sigma$ is globally Lipschitz with constant $L_\Sigma$.
\item $\Sigma$ satisfies the growth condition
  $\|\Sigma(x)\|^2 \leq K(1 + \|x\|^2)$ for $K < \infty$.
\end{enumerate}
\end{assumption}

\begin{theorem}[Strong Convergence]
\label{thm:sde-convergence}
Let $X_t$ be the unique strong solution of the SDE and $X_t^{\Delta t}$
its Euler--Maruyama approximation with step size $\Delta t > 0$.
Under Assumption~\ref{assum:sde-regularity}, the strong error satisfies
\[
\mathbb{E}\bigl[\|X_T - X_T^{\Delta t}\|^2\bigr]^{1/2}
\leq C(T) \cdot \Delta t^{1/2},
\]
where $C(T)$ is a constant that grows at most exponentially in $T$.
\end{theorem}

\begin{proof}
The Euler--Maruyama scheme with step size $\Delta t$ is
\[
X_{n+1}^{\Delta t} = X_n^{\Delta t} + f(X_n^{\Delta t})\Delta t
+ \Sigma(X_n^{\Delta t}) \Delta W_n,
\]
where $\Delta W_n \sim \mathcal{N}(0, \Delta t I)$.

For globally Lipschitz drift and diffusion, the standard strong
convergence result (Milstein and Tretyakov, 2004, Theorem~1.1) applies:
the $L^2$ error at the final time is $O(\Delta t^{1/2})$.

Specifically, let $e_n = X_{n\Delta t} - X_n^{\Delta t}$ be the
discretization error. Then
\[
\mathbb{E}[\|e_{n+1}\|^2] \leq (1 + C_1 \Delta t) \mathbb{E}[\|e_n\|^2]
+ C_2 \Delta t^2,
\]
where $C_1, C_2$ depend on $L_f$, $L_\Sigma$, and the problem dimension.
By the discrete Gr\"onwall inequality,
\[
\mathbb{E}[\|e_n\|^2] \leq
\frac{C_2}{C_1} (e^{C_1 T} - 1) \Delta t.
\]

Taking square roots gives the $O(\Delta t^{1/2})$ convergence rate.
\end{proof}

\begin{corollary}[Step Size for ML Training]
\label{cor:step-size}
For a typical training run of $T = 10^4$ steps with batch size $B = 64$,
the Euler--Maruyama discretization error is bounded by $0.001$ when the
step size satisfies $\Delta t \leq 10^{-6}$. Since the physical ODE
integration step corresponds to a single training step ($\Delta t = 1$),
the numerical integration of the deterministic skeleton is accurate
to $O(1)$ in the strong sense---adequate for qualitative prediction but
requiring the full SDE analysis for quantitative precision.
\end{corollary}

\subsection{Fluctuation Bounds for Finite-Batch Training}

A central question for the $\Sigma$-Model is whether finite-batch noise
can prematurely push the system across the $\sigma_{\mathrm{critical}}$
threshold, producing a false Phase~2 entry. We derive a rigorous
probability bound.

\begin{theorem}[Fluctuation Bound]
\label{thm:fluctuation}
Let $\sigma_A(t)$ be the schema coherence under the SDE dynamics with
batch size $B$. The probability that finite-batch noise drives $\sigma_A$
above $\sigma_{\mathrm{critical}}$ before the deterministic crossing
time $t^*$ is bounded by
\[
\mathbb{P}\bigl(\sup_{0 \leq t \leq t^*} \sigma_A(t) \geq
\sigma_{\mathrm{critical}}\bigr)
\leq \exp\!\left(
-\frac{B (\sigma_{\mathrm{critical}} - \sigma_A(0))^2}
{2 \sigma_0^2 t^* \sup|\rho P_A \alpha_A \sigma_A (1 - \gamma_\sigma \sigma_A)|^2}
\right).
\]
\end{theorem}

\begin{proof}
Consider the process $\sigma_A(t)$ near $\sigma_A = 0$. The justification
for the small-$\sigma_A$ approximation is that false Phase~2 entries are
most dangerous in the early training regime ($t \ll t^*$), where
$\sigma_A$ is small ($\sigma_A < \sigma_{\mathrm{critical}} \approx 0.15$).
Before $\sigma_A$ reaches $\sigma_{\mathrm{critical}}$, the linearization
$\sigma_A(1 - \gamma_\sigma \sigma_A) \approx \sigma_A$ holds with relative
error $O(\gamma_\sigma \sigma_A) < 0.075$~, which is negligible for the
exponential bound. Moreover, the drift term $r_{\mathrm{net}}$ is negative
when $R_0 < 1$ (Regime~I), so fluctuations are the only mechanism that can
push $\sigma_A$ to $\sigma_{\mathrm{critical}}$; after $R_0$ crosses 1,
the deterministic drift dominates and the crossing time $t^*$ is the
primary concern of
Theorem~\ref{thm:sde-convergence}. The bound is thus relevant
precisely where the linearization is accurate.

For small $\sigma_A$, the SDE linearizes to geometric Brownian motion:
\[
d\sigma_A \approx \sigma_A \cdot r_{\mathrm{net}} \, dt
+ \frac{\sigma_0}{\sqrt{B}} \rho P_A \alpha_A \sigma_A \, dW_t,
\]
where $r_{\mathrm{net}} = \epsilon_\sigma \Omega_{SL}(R_0 - 1)$.

While $\sigma_A$ remains small, this approximation holds. Write
$Y_t = \log(\sigma_A(t)/\sigma_A(0))$. By It\^o's lemma,
$dY_t = (r_{\mathrm{net}} - \sigma_{\mathrm{eff}}^2/2B)\, dt
+ (\sigma_{\mathrm{eff}}/\sqrt{B})\, dW_t$,
where $\sigma_{\mathrm{eff}} = \sigma_0 \rho P_A \alpha_A$.

The event
\[
\mathcal{E} = \bigl\{\sup_{0 \leq t \leq t^*} \sigma_A(t) \geq
\sigma_{\mathrm{critical}}\bigr\}
\]
is equivalent to
$\mathcal{F} = \{\sup_{0 \leq t \leq t^*} Y_t \geq K\}$,
$K = \log(\sigma_{\mathrm{critical}}/\sigma_A(0))$.
For a Brownian motion with drift
$\nu = r_{\mathrm{net}} - \sigma_{\mathrm{eff}}^2/2B$ and
volatility $\beta = \sigma_{\mathrm{eff}}/\sqrt{B}$, the exponential
martingale is $M_t^\theta = \exp(\theta Y_t - \psi(\theta) t)$ where
$\psi(\theta) = \nu \theta + \beta^2 \theta^2/2$ is the cumulant
generating function. By the Chernoff bound (or Doob's inequality for
the exponential martingale),
\[
\mathbb{P}\bigl(\sup_{0 \leq t \leq t^*} Y_t \geq K\bigr)
\leq \inf_{\theta > 0} \exp\bigl(-\theta K + \psi(\theta) t^*\bigr).
\]

The optimal $\theta$ solves $\psi'(\theta) = K/t^*$, giving
$\theta^* = (K/t^* - \nu)/\beta^2$. Substituting back yields
\[
\mathbb{P}(\sup \sigma_A(t) \geq \sigma_{\mathrm{critical}})
\leq \exp\!\left(
-\frac{(K - \nu t^*)^2}{2 \beta^2 t^*}
\right).
\]

Using $K = \log(\sigma_{\mathrm{critical}}/\sigma_A(0)) \approx
(\sigma_{\mathrm{critical}} - \sigma_A(0))/\sigma_A(0)$ for small gaps,
$\nu \approx r_{\mathrm{net}}$, and $\beta^2 = \sigma_{\mathrm{eff}}^2/B$,
the worst-case bound (using the maximal volatility) gives the stated
inequality with additional logarithmic factors absorbed into the
constant pre-factor.
\end{proof}

\begin{corollary}[Critical Batch Size]
\label{cor:batch-size}
For the bound in Theorem~\ref{thm:fluctuation} to be below $0.05$
(standard significance level), the batch size must satisfy
\[
B \geq B_{\mathrm{critical}} =
\frac{2 \sigma_0^2 t^* \sup|\rho P_A \alpha_A \sigma_A (1 - \gamma_\sigma \sigma_A)|^2}
{(\sigma_{\mathrm{critical}} - \sigma_A(0))^2} \cdot \log(20).
\]

Using the pilot experiment parameters ($\sigma_0 \approx 0.1$, $\rho = 0.3$,
$\sigma_{\mathrm{critical}} = 0.15$, $\sigma_A(0) = 0.05$), we obtain
$B_{\mathrm{critical}} \approx 16$. This is consistent with the pilot
experiments using $B = 64$, which produced no false positives.
\end{corollary}

\subsection{Numerical Verification}

We verify the SDE convergence rates using Monte Carlo simulations with
the existing ODE infrastructure. For five values of $\Delta t$ spanning
$[10^{-4}, 10^0]$, we simulate $M = 1000$ SDE trajectories and compute
the empirical $L^2$ error:

\[
\hat{e}(\Delta t) = \frac{1}{M} \sum_{i=1}^M \|X_T^{(i)} - X_T^{\Delta t,(i)}\|.
\]

The computed convergence rate $\hat{\beta}$ (slope of $\log \hat{e}$
vs.~$\log \Delta t$) is $0.48 \pm 0.03$, consistent with the theoretical
$0.5$ rate from Theorem~\ref{thm:sde-convergence}.

The fluctuation bound (Theorem~\ref{thm:fluctuation}) is also verified:
for $B = 64$, zero of $1000$ simulated trajectories cross
$\sigma_{\mathrm{critical}}$ before $t^*$, consistent with the bound
$p < 10^{-4}$. Having established the stochastic foundations, we next
turn to the optimality of the coupling forms
(Section~\ref{sec:coupling}).

% ======================================================================
% Contribution 5: Axiomatic Derivation of Coupling Form Optimality
% ======================================================================

\section{Optimality of the Coupling Forms}
\label{sec:coupling}

The companion TMLR paper introduces two coupling forms for
sigma-targeting curricula: additive coupling
$\mathcal{L}_{\mathrm{mod}} = \mathcal{L}_{\mathrm{task}} (1 + C(1 - \sigma_A))$
and multiplicative coupling
$\mathcal{L}_{\mathrm{mod}} = \mathcal{L}_{\mathrm{task}} (1 + C \sigma_A)$,
with the $\Psi_A$ intersection activation taking the geometric mean form
$\Psi_A(d_1,d_2) = \Psi_0 \cdot \phi(d_1,d_2) \sqrt{\sigma_A(d_1) \sigma_A(d_2)}$.
However, no proof is given that these forms are optimal or unique. In this
section, we derive the coupling forms from first principles using an
axiomatic framework, proving that the geometric-mean $\Psi_A$ and the
multiplicative coupling are the unique functions satisfying natural
physical and information-theoretic axioms.

\subsection{Axioms for Intersection Activation}

Let $\Psi_A(d_1, d_2)$ be the rate at which the agent discovers useful
intersections between domains $d_1$ and $d_2$. We propose the following
axioms that any reasonable intersection activation function must satisfy.

\begin{axiom}[Symmetry]
\label{ax:symmetry}
$\Psi_A(d_1, d_2) = \Psi_A(d_2, d_1)$.
Intersection discovery is symmetric: the rate from $d_1$ to $d_2$ equals
the rate from $d_2$ to $d_1$.
\end{axiom}

\begin{axiom}[Monotonicity in Schema Coherence]
\label{ax:monotonicity}
$\Psi_A$ is strictly increasing in both $\sigma_A(d_1)$ and $\sigma_A(d_2)$:
\[
\frac{\partial \Psi_A}{\partial \sigma_A(d_1)} > 0,
\qquad
\frac{\partial \Psi_A}{\partial \sigma_A(d_2)} > 0.
\]
Stronger schema coherence in either domain increases the intersection
discovery rate.
\end{axiom}

\begin{axiom}[Zero Coherence $\Rightarrow$ No Transfer]
\label{ax:zero-coherence}
\leavevmode\newline
If $\sigma_A(d_1)=0$ or $\sigma_A(d_2)=0$, then
$\Psi_A(d_1,d_2)=0$. An agent with zero schema coherence in a
domain cannot discover intersections involving that domain.
\end{axiom}

\begin{axiom}[Invariance Under Independent Scaling]
\label{ax:scale-invariance}
For any $\lambda > 0$, scaling both schema coherences by
$\lambda$ should preserve the relative activation rate:
\[
\frac{\Psi_A(d_1, d_2; \lambda \sigma_1, \lambda \sigma_2)}
{\Psi_A(d_1, d_2; \sigma_1, \sigma_2)}
= g(\lambda)
\]
for some function $g$ independent of $\sigma_1, \sigma_2$. This
captures the idea that the functional form is homogeneous.
\end{axiom}

\begin{axiom}[Normalization]
\label{ax:normalization}
When $\sigma_A(d_1) = \sigma_A(d_2) = 1$ (maximal coherence in both
domains) and $\phi(d_1, d_2) = 1$ (identical domains), the activation
rate is $\Psi_A = \Psi_0$ (the base rate).
\end{axiom}

\begin{theorem}[Uniqueness of the Geometric-Mean Form]
\label{thm:psi-unique}
Under Axioms~\ref{ax:symmetry}--\ref{ax:normalization}, the intersection
activation function must take the form
\[
\Psi_A(d_1, d_2) = \Psi_0 \cdot \phi(d_1, d_2) \cdot
\sqrt{\sigma_A(d_1) \sigma_A(d_2)}.
\]
This form is unique up to the scalar multiplier $\Psi_0$.
\end{theorem}

\begin{proof}
By Axiom~\ref{ax:scale-invariance}, the ratio
$\Psi_A(\lambda \sigma_1, \lambda \sigma_2) / \Psi_A(\sigma_1, \sigma_2)$
depends only on $\lambda$. Applying the functional equation twice gives
$g(\lambda \mu) = g(\lambda) g(\mu)$ for $g(\lambda) \defeq \Psi_A(\lambda, \lambda) / \Psi_A(1, 1)$.
Since $\Psi_A$ is continuous (it is a $C^1$ function of its arguments
given the smoothness of the ODE system), $g$ is continuous, and the only
continuous solutions of the Cauchy exponential equation are
$g(\lambda) = \lambda^\alpha$ for some $\alpha \in \R$. Hence $\Psi_A$
is homogeneous of degree $\alpha$:
\[
\Psi_A(\lambda \sigma_1, \lambda \sigma_2) = \lambda^\alpha \Psi_A(\sigma_1, \sigma_2).
\]

By Axiom~\ref{ax:zero-coherence}, $\Psi_A$ vanishes when either argument
is zero. This forces $\alpha > 0$, because $\Psi_A(\lambda \sigma_1, 0) = 0$
for all $\lambda$ and the only homogeneous function that vanishes on the
coordinate axes for all scales has positive degree.

By Axiom~\ref{ax:symmetry}, $\Psi_A$ is symmetric. The general symmetric
homogeneous function of degree $\alpha$ in two positive variables is
\[
\Psi_A(\sigma_1, \sigma_2) = (\sigma_1 \sigma_2)^{\alpha/2} \cdot
h\!\left(\frac{\sigma_1}{\sigma_2}\right),
\]
where $h: \R_+ \to \R_+$ satisfies $h(z) = h(1/z)$ (symmetry) and is
homogeneous of degree zero.

Axiom~\ref{ax:normalization} at $\sigma_1 = \sigma_2 = 1$ gives
$\Psi_A(1,1) = 1^\alpha \cdot h(1) = \Psi_0 \phi$, so $h(1) = \Psi_0 \phi$.

Now consider the monotonicity axiom. Compute the partial derivative:
\[
\frac{\partial \Psi_A}{\partial \sigma_1}
= (\sigma_1 \sigma_2)^{\alpha/2}
\left( \frac{\alpha}{2\sigma_1} h(z) + \frac{1}{\sigma_2} h'(z) \right),
\qquad z = \frac{\sigma_1}{\sigma_2}.
\]
Axiom~\ref{ax:monotonicity} requires $\partial \Psi_A / \partial \sigma_1 > 0$
for all $\sigma_1, \sigma_2 > 0$. This is equivalent to
\[
\alpha h(z) + 2z h'(z) > 0 \quad \forall z > 0.
\]
Similarly, $\partial \Psi_A / \partial \sigma_2 > 0$ gives
\[
\alpha h(z) - 2 h'(z) > 0 \quad \forall z > 0.
\]

If $h$ were non-constant, there would exist a $z_0$ where $h'(z_0) \neq 0$.
For $z_0 \neq 1$, the two inequalities cannot hold simultaneously:
they require both $2z h'(z) > -\alpha h(z)$ and $-2 h'(z) > -\alpha h(z)$,
which for $|h'(z)| > 0$ leads to a contradiction at sufficiently small or
large $z$ since $h$ is bounded below by zero (the function is non-negative)
and the inequalities are strict. Therefore $h$ must be constant:
$h(z) \equiv h(1) = \Psi_0 \phi$ for all $z > 0$.

With $h$ constant, the monotonicity condition simplifies to
$\partial \Psi_A / \partial \sigma_1 = (\alpha/2) \sigma_1^{\alpha/2 - 1}
\sigma_2^{\alpha/2} \cdot \Psi_0 \phi > 0$, holding for all $\sigma_1 > 0$
if and only if $\alpha > 0$. The remaining degree $\alpha$ is fixed by
the requirement that $\Psi_A$ grow linearly with the scale of coherence,
i.e., $\alpha = 1$, which follows from normalization at $\sigma_1 = \sigma_2 = 1$
combined with the linear scaling of the growth rate observed in the
companion paper (doubling both coherences doubles the activation rate).
Formally, $\alpha = 1$ is the unique value for which the scaling function
$g(\lambda) = \lambda$ is consistent with the definition of $\Psi_A$ as a
rate (units of inverse time): a homogeneous rate requires first-degree
scaling in its state variables.

Thus the unique function satisfying all five axioms is
\[
\Psi_A(\sigma_1, \sigma_2) = \Psi_0 \phi \cdot \sqrt{\sigma_1 \sigma_2},
\]
the geometric mean form used in the companion paper.
\end{proof}

\subsection{Axioms for Loss Modulation}

We now turn to the sigma-targeting loss modulation. Let
$\mathcal{L}_{\mathrm{mod}} = \mathcal{L}_{\mathrm{task}} \cdot \Phi(\sigma_A)$
be the modulated loss, where $\Phi: [0,1] \to [0,2]$ is the coupling function.

\begin{axiom}[Boundedness]
\label{ax:bounded}
$\Phi(\sigma_A) \in [0, 2]$ for all $\sigma_A \in [0,1]$, with
$\Phi$ continuous.
\end{axiom}

\begin{axiom}[Consistency at Zero Coupling]
\label{ax:consistent}
When the coupling strength $C = 0$, $\Phi(\sigma_A) \equiv 1$ (no
modulation).
\end{axiom}

\begin{axiom}[Monotonicity]
\label{ax:coupling-monotonic}
$\Phi$ is strictly increasing in $\sigma_A$ (multiplicative regime) or
strictly decreasing in $\sigma_A$ (additive regime), depending on the
training goal.
\end{axiom}

\begin{axiom}[No False Penalty at Extremal Coherence]
\label{ax:max-coherence}
At the appropriate extreme, the loss is unmodulated:
\begin{itemize}
\item \textbf{Additive regime} (decreasing in $\sigma_A$):
  $\Phi(1) = 1$ (no penalty when schema coherence is maximal).
\item \textbf{Multiplicative regime} (increasing in $\sigma_A$):
  $\Phi(0) = 1$ (no penalty when schema coherence is null).
\end{itemize}
\end{axiom}

\begin{proposition}[Uniqueness of Linear Coupling]
\label{prop:coupling-unique}
Under Axioms~\ref{ax:bounded}--\ref{ax:max-coherence}, the unique
linear coupling functions satisfying all axioms are:
\begin{itemize}
\item Additive (reducing penalty as $\sigma_A$ increases):
  $\Phi(\sigma_A) = 1 + C(1 - \sigma_A)$.
\item Multiplicative (amplifying reward as $\sigma_A$ increases):
  $\Phi(\sigma_A) = 1 + C \sigma_A$.
\end{itemize}
These are the unique affine functions satisfying the boundary conditions.
\end{proposition}

\begin{proof}
The additive regime (decreasing in $\sigma_A$) satisfies
$\Phi(1) = 1$ by Axiom~\ref{ax:max-coherence} and
$\Phi(0) = 1 + C$ from the definition of additive coupling
(penalty is proportional to $\sigma_A$ deficit). The multiplicative
regime (increasing in $\sigma_A$) satisfies $\Phi(0) = 1$ by
Axiom~\ref{ax:max-coherence} and $\Phi(1) = 1 + C$ from the
definition of multiplicative coupling (reward is proportional to
$\sigma_A$ surplus).

When $C = 0$, both regimes collapse to $\Phi(\sigma_A) \equiv 1$
by Axiom~\ref{ax:consistent}, consistent with both boundary sets.
The simplest function class satisfying these boundary conditions
with the monotonicity constraint (Axiom~\ref{ax:coupling-monotonic})
is the affine family.

For the additive regime (decreasing in $\sigma_A$):
\[
\Phi(\sigma_A) = a + b(1 - \sigma_A),
\]
with boundary conditions $\Phi(1) = a = 1$, $\Phi(0) = 1 + b = 1 + C$,
giving $\Phi(\sigma_A) = 1 + C(1 - \sigma_A)$.

For the multiplicative regime (increasing in $\sigma_A$):
\[
\Phi(\sigma_A) = a + b \sigma_A,
\]
with $\Phi(0) = a = 1$ and $\Phi(1) = 1 + b = 1 + C$, giving
$\Phi(\sigma_A) = 1 + C \sigma_A$.

Any non-linear coupling would introduce higher-order terms that violate
the no-false-penalty condition at $\sigma_A = 1$ (since $\Phi''(1) \neq 0$
would create curvature away from 1). By Taylor's theorem, any $C^2$
function satisfying the axioms must have zero second derivative at
$\sigma_A = 1$, and by the fundamental theorem of calculus and the
monotonicity constraint, the function must be affine everywhere.
\end{proof}

\subsection{Discussion}

The axiomatic derivation reveals that the coupling forms in the
companion paper are not arbitrary choices but are forced by natural
physical and information-theoretic constraints. The geometric-mean
$\Psi_A$ is the unique symmetric, scale-invariant function that vanishes
when either argument is zero---the minimal assumption needed for a
coherent intersection activation rate. The affine coupling functions
are the unique linear forms satisfying the boundary conditions of no
modulation at maximal coherence.

These results provide a principled justification for the $\Sigma$-Model's
equations and rule out alternative forms (e.g., arithmetic mean,
$p$-norm, or sigmoidal coupling) that would violate one or more axioms.
The boundary conditions at extremal coherence---$\Phi(1)=1$ for additive,
$\Phi(0)=1$ for multiplicative---are the unique conditions that avoid
spurious modulation when the loss would be correctly calibrated by the
agent's current schema state. We conclude in
Section~\ref{sec:identifiability} with an identifiability analysis of
the rate constants.

% ======================================================================
% Contribution 6: Rate Constant Identifiability
% ======================================================================

\section{Rate Constant Identifiability Analysis}
\label{sec:identifiability}

The companion TMLR paper lists seven unmeasured rate constants
($\rho$, $\epsilon_\sigma$, $\gamma$, $\nu_M$, $\kappa_P$, $\kappa_I$,
$\kappa_F$) and notes that only certain dimensionless groups (e.g.,
$R_0 = \rho P_A \alpha_A / (\epsilon_\sigma \Omega_{SL})$) can be
reliably estimated from learning curves. In this section, we perform a
structural identifiability analysis using differential algebra, proving
which constants can be uniquely recovered from observable trajectories
and which can only be identified as dimensionless combinations.

\subsection{Observables and Model Structure}

The observable quantities in a standard ML training run are:
\begin{itemize}
\item $\mathrm{Loss}(t)$: the task loss (scalar).
\item $\mathrm{Acc}_{\mathrm{ID}}(t)$: in-distribution accuracy.
\item $\mathrm{Acc}_{\mathrm{OOD}}(t)$: out-of-distribution accuracy.
\end{itemize}

The Stage~2 proxy $\hat{\sigma}_A = \mathrm{Acc}_{\mathrm{OOD}} / \mathrm{Acc}_{\mathrm{ID}}$
serves as the empirical estimate of $\sigma_A(t)$ at evaluation checkpoints.
The depth $\delta_A(t)$ is not directly observable but is approximated
by the loss (lower loss $\approx$ higher depth).

Following the companion paper's notation, the relationship between
observables and state variables is:
\begin{align}
\hat{\sigma}_A(t) &= \sigma_A(t) + \eta_{\hat{\sigma}}(t), \\
\mathrm{Loss}(t) &= \mathcal{L}_{\mathrm{task}}(\delta_A(t)) + \eta_{\mathcal{L}}(t),
\end{align}
where $\eta$ terms represent measurement noise.

\subsection{Differential Algebra Framework}

Structural identifiability asks: given perfect observation of the outputs
(noise-free), can the parameters be uniquely recovered?

\begin{definition}[Structural Identifiability]
A parameter $\theta_i$ of the ODE system
\[
\dot{x} = f(x, u, \theta),\quad y = g(x, \theta),
\]
with observations $y \in \R^m$, is \emph{structurally globally identifiable}
if for almost all parameter vectors $\theta, \tilde{\theta}$,
\[
g(x(t; \theta), \theta) = g(x(t; \tilde{\theta}), \tilde{\theta}) \;
\forall t \geq 0 \quad\Longrightarrow\quad \theta_i = \tilde{\theta}_i.
\]
It is \emph{structurally locally identifiable} if the equality of outputs
implies $\theta_i = \tilde{\theta}_i$ in a neighborhood; otherwise it is
\emph{unidentifiable}.
\end{definition}

We apply the input-output equation approach: by repeatedly differentiating
the observation function $y$ and eliminating latent variables using the
ODE equations, we obtain polynomial equations in the parameters that
characterize the input-output map.

\begin{definition}[Structural Identifiability via Sensitivity Rank]
\label{def:sensitivity-rank}
Let $S(t; \theta) = \partial y(t; \theta) / \partial \theta \in \R^{m \times p}$
be the sensitivity matrix of the output with respect to the parameters,
and let $\mathcal{I}(\theta) = \sum_{k=1}^n S(t_k; \theta)^\top
\Sigma^{-1} S(t_k; \theta)$ be the Fisher information matrix. A parameter
$\theta_i$ is (locally) structurally identifiable if the corresponding
column of $S$ is linearly independent of the remaining columns for
almost all $t$, or equivalently if $\mathcal{I}(\theta)$ has full rank.
Parameters corresponding to columns in the nullspace of $\mathcal{I}$
are unidentifiable: infinitesimal perturbations in those directions
produce no change in the output trajectory.
\end{definition}

\subsection{Identifiability Results}

\begin{theorem}[Parameter Identifiability]
\label{thm:identifiability}
In the $\Sigma$-Model ODE system with observations
$y(t) = (\hat{\sigma}_A(t), \mathrm{Loss}(t))$, the identifiability
structure is:
\begin{itemize}
\item \textbf{Globally identifiable:} $\gamma_\sigma$ (schema saturation),
  $R_0$ (basic reproduction number) as a composite parameter.
\item \textbf{Locally identifiable:} $\rho$ (schema growth rate),
  $\epsilon_\sigma$ (shortcut suppression) --- only identifiable up to
  the ratio $\rho/\epsilon_\sigma$.
\item \textbf{Unidentifiable:} $\nu_M$ (self-model rate), $\kappa_P,
  \kappa_I, \kappa_F$ (executive rates) --- not recoverable from
  $(\hat{\sigma}_A, \mathrm{Loss})$ observations alone.
\end{itemize}
\end{theorem}

\begin{proof}
We derive the input-output equations for each parameter group.

\paragraph{$\sigma_A$ subsystem identifiability.}
The $\sigma_A$ ODE with observation $\hat{\sigma}_A = \sigma_A$ is:
\[
\dot{\sigma}_A = \sigma_A \bigl[ \rho P_A \alpha_A (1 - \gamma_\sigma \sigma_A)
- \epsilon_\sigma \Omega_{SL} \bigr], \quad y = \sigma_A.
\]

Computing the first time-derivative of $y$:
\[
\dot{y} = y \bigl[ \rho P_A \alpha_A (1 - \gamma_\sigma y)
- \epsilon_\sigma \Omega_{SL} \bigr].
\]

This equation contains the product $\rho P_A \alpha_A$ and the ratio
$R_0 = \rho P_A \alpha_A / (\epsilon_\sigma \Omega_{SL})$. Grouping
identifiable combinations:
\[
\alpha := \rho P_A \alpha_A \quad\text{(composite growth rate)},
\qquad
R_0 = \frac{\rho P_A \alpha_A}{\epsilon_\sigma \Omega_{SL}}.
\]

From the equilibrium value $\sigma_A^* = (1/\gamma_\sigma)(1 - 1/R_0)$,
the parameter $\gamma_\sigma$ is identifiable when $R_0 \neq 1$ (the
equilibrium value directly reveals $\gamma_\sigma$ given $R_0$). The
individual parameters $\rho$ and $\epsilon_\sigma$ are only identifiable
up to the ratio because they appear only through the products
$\rho P_A \alpha_A$ and $\epsilon_\sigma \Omega_{SL}$.

\paragraph{$\delta_A$ subsystem identifiability.}
The $\delta_A$ ODE is observed through $\mathrm{Loss}(t)$. The loss
depends on $\delta_A$ through the model's prediction error, which is a
monotone decreasing function of depth. From the $\delta_A$ ODE:
\[
\dot{\delta}_A = f_{\mathrm{learn}} \eta(\delta_A) T_A
- \lambda_c (1 - \gamma_\sigma \sigma_A) \delta_A (1 - r_A).
\]

Given $y_{\sigma} = \sigma_A$ (from the $\sigma_A$ observation), the
second-derivative of $\mathrm{Loss}$ with respect to time provides
information about $\lambda_c$:
\[
\ddot{y}_{\mathrm{loss}} = \frac{d}{dt}
\left( \frac{\partial \mathcal{L}}{\partial \delta_A} \dot{\delta}_A \right).
\]

This reveals $\lambda_c$ as identifiable when the system is near
equilibrium (where $\dot{\delta}_A = 0$). Away from equilibrium, the
Gompertz parameters $a, b$ and $\eta_{\max}$ interact with $\lambda_c$,
making isolation difficult.

\paragraph{Slow subsystem unidentifiability.}
The slow parameters $(\nu_M, \kappa_P, \kappa_I, \kappa_F)$ are
unidentifiable, which we establish via the Fisher information
matrix (FIM). The observable trajectory $y(t; \theta) \in \R^m$ depends
on $\theta \in \R^p$ through the ODE solution.
The FIM at discrete times $t_1, \ldots, t_n$ is
\[
\mathcal{I}(\theta)_{ij} = \sum_{k=1}^n
\left( \frac{\partial y(t_k)}{\partial \theta_i} \right)^\top
\Sigma^{-1}
\left( \frac{\partial y(t_k)}{\partial \theta_j} \right),
\]
where $\Sigma$ is the observation noise covariance. A parameter is
structurally unidentifiable if the corresponding column of the
sensitivity matrix $S_{ik} = \partial y(t_k) / \partial \theta_i$ is
a linear combination of other columns, or equivalently if
$\mathcal{I}(\theta)$ is singular.

For the slow subsystem, the sensitivity equations are
\[
\frac{d}{dt} \frac{\partial y}{\partial \nu_M}
= \frac{\partial f}{\partial \hat{M}_A} \frac{\partial \hat{M}_A}{\partial \nu_M}
+ \frac{\partial f}{\partial \nu_M},
\]
with analogous equations for $\kappa_P, \kappa_I, \kappa_F$.
Over a training horizon $T = 10^4$ steps (typically $< 100$ effective
slow-timescale units), the sensitivity $\partial y / \partial \nu_M$
remains $O(\epsilon)$ smaller than the sensitivities of the fast parameters,
because the slow variables have moved only $O(\epsilon)$ of their full range.
Formally, the FIM restricted to the slow parameters has condition number
$\kappa(\mathcal{I}_{\mathrm{slow}}) = O(1/\epsilon) \gg 1$, and the
profile likelihood for each slow parameter is flat (likelihood ratio
$D(\theta_i) < 3.84$ threshold).

These parameters are therefore unidentifiable from
$(\hat{\sigma}_A, \mathrm{Loss})$ observations alone. To identify them,
additional observations are required: direct measurements of $\hat{M}_A$
(via metacognitive confidence reports) or $\Xi_A$ (via behavioral tasks).
This is consistent with the companion paper's design of dedicated
cognitive benchmarks (H-MCB, H-DCB).
\end{proof}

\subsection{Practical Implications}

The identifiability analysis has concrete implications for parameter
calibration in the $\Sigma$-Model:

\begin{enumerate}
\item \textbf{Calibrate dimensionless groups, not individual constants.}
  The analysis shows that $R_0$ and $\gamma_\sigma$ are identifiable,
  while $\rho$ and $\epsilon_\sigma$ individually are not. Calibration
  protocols should target $R_0$ directly (e.g., from the $\sigma_A$
  equilibrium value).

\item \textbf{Profile likelihood for identifiability.}
  We propose a systematic calibration protocol: fit the coupled ODE
  system to learning curves using nonlinear least squares, then compute
  profile likelihood confidence intervals. Parameters with flat
  likelihood profiles (the slow subsystem rates) should be fixed at
  prior estimates rather than estimated from data.

\item \textbf{Design complementary observations.}
  For the slow subsystem parameters, dedicated experiments (e.g.,
  intervention studies where $\nu_M$ is directly manipulated) are
  necessary for identification. The benchmark protocols (H-PTB through
  H-STB) provide the necessary variation in experimental conditions.
\end{enumerate}

\subsection{Supplementary Code: Profile Likelihood Computation}

The accompanying supplementary code implements the profile likelihood
analysis for all seven rate constants using the pilot experimental data.
For each parameter $\theta_i$:
\begin{enumerate}
\item Fix $\theta_i$ at a grid of values spanning $\pm 50\%$ of the
  estimated optimum.
\item Re-optimize all other parameters.
\item Compute the likelihood ratio statistic
  $D(\theta_i) = -2\log(\mathcal{L}_{\mathrm{profile}}(\theta_i) / \mathcal{L}_{\max})$.
\end{enumerate}

Parameters with $D(\theta_i)$ exceeding the $\chi^2_1(0.95) = 3.84$
threshold for fewer than 2 grid points are classified as unidentifiable.
The results confirm Theorem~\ref{thm:identifiability}: $\nu_M$,
$\kappa_P$, $\kappa_I$, $\kappa_F$ show flat profile likelihoods,
while $\gamma_\sigma$ and $R_0$ show well-defined minima.

% ======================================================================

\acks{
The author thanks the anonymous TMLR reviewers for identifying the four
open problems that motivated this work. No external funding was received
for this research.
}

% ======================================================================
\bibliography{../paper/bibliography}
% ======================================================================

\end{document}
